\documentclass[11pt]{article}

\usepackage[T1]{fontenc}
\usepackage[utf8]{inputenc}
\usepackage[english]{babel}
\usepackage[margin=1in]{geometry}
\usepackage{lmodern}
\usepackage{microtype}
\usepackage{mathtools}
\usepackage{amsmath,amssymb,amsthm}
\usepackage{array}
\usepackage{booktabs}
\usepackage{tabularx}
\usepackage{float}
\usepackage{graphicx}
\usepackage{caption}
\usepackage{enumitem}
\usepackage{fancyhdr}
\usepackage[numbers,sort&compress]{natbib}
\usepackage{doi}
\usepackage{xurl}
\usepackage{hyperref}
\usepackage{cleveref}
\usepackage{orcidlink}
\usepackage{titling}



\DeclareMathOperator{\Capacity}{Cap}
\DeclareMathOperator{\Cost}{Cost}
\DeclareMathOperator{\Ran}{Ran}
\DeclareMathOperator{\Ker}{Ker}
\DeclareMathOperator{\tr}{tr}
\DeclareMathOperator{\dist}{dist}

\newcommand{\K}{\mathsf K}
\newcommand{\XiRes}{\Xi}
\newcommand{\preceqq}{\preceq}
\newcommand{\eps}{\varepsilon}

\newcommand{\DBG}{D^{\mathrm{BG}}}
\newcommand{\LBG}{L^{\mathrm{BG}}}
\newcommand{\Lphys}{L^{\mathrm{phys}}}
\newcommand{\Tmax}{T_{\max}}
\newcommand{\Omsrc}{\Omega_{\mathrm{src}}}
\newcommand{\Omtr}{\Omega_{\mathrm{tr}}}
\newcommand{\Ompr}{\Omega_{\mathrm{pr}}}
\newcommand{\Omvisc}{\Omega_{\mathrm{visc}}}
\newcommand{\Omtail}{\Omega_{\mathrm{tail}}}
\newcommand{\Omstr}{\Omega_{\mathrm{str}}}
\newcommand{\Omamp}{\Omega_{\mathrm{amp}}}
\newcommand{\Omend}{\Omega_{\mathrm{endpoint}}}
\newcommand{\CG}{C_G}
\newcommand{\sigmaJ}{\sigma_J}
\newcommand{\lambdaJ}{\lambda_J}
\newcommand{\LPtr}{S_j}
 
\newtheorem{theorem}{Theorem}[section]
\newtheorem{lemma}[theorem]{Lemma}
\newtheorem{proposition}[theorem]{Proposition}
\newtheorem{corollary}[theorem]{Corollary}
\theoremstyle{definition}
\newtheorem{definition}[theorem]{Definition}
\newtheorem{remark}[theorem]{Remark}
\newtheorem{nonclaim}[theorem]{Nonclaim}
\floatstyle{ruled}
\newfloat{algorithm}{tbp}{loa}
\floatname{algorithm}{Algorithm}

\urlstyle{same}
\captionsetup{
  font=small,
  labelfont=bf,
  labelsep=period
}
\captionsetup[table]{skip=6pt,position=top}
\captionsetup[figure]{skip=6pt}
\captionsetup[algorithm]{skip=6pt,position=top}
\crefname{algorithm}{algorithm}{algorithms}
\Crefname{algorithm}{Algorithm}{Algorithms}
\crefname{theorem}{theorem}{theorems}
\Crefname{theorem}{Theorem}{Theorems}
\crefname{lemma}{lemma}{lemmas}
\Crefname{lemma}{Lemma}{Lemmas}
\crefname{proposition}{proposition}{propositions}
\Crefname{proposition}{Proposition}{Propositions}
\crefname{corollary}{corollary}{corollaries}
\Crefname{corollary}{Corollary}{Corollaries}
\crefname{definition}{definition}{definitions}
\Crefname{definition}{Definition}{Definitions}
\crefname{remark}{remark}{remarks}
\Crefname{remark}{Remark}{Remarks}
\setlist[enumerate]{leftmargin=*, itemsep=2pt, topsep=4pt}
\setlength{\droptitle}{-3.2em}
\pretitle{\begin{center}\LARGE\bfseries}
\posttitle{\par\end{center}\vspace{0.5em}}
\preauthor{\begin{center}\normalsize}
\postauthor{\par\end{center}\vspace{-0.4em}}
\predate{\begin{center}\small}
\postdate{\par\end{center}\vspace{0.6em}}
\setlength{\emergencystretch}{2em}
\clubpenalty=10000
\widowpenalty=10000
\displaywidowpenalty=10000
\interfootnotelinepenalty=10000
\allowdisplaybreaks[2]
\fancypagestyle{firstpage}{\fancyhf{}
  \fancyfoot[C]{\scriptsize Automorph Inc., Wilmington, DE, USA \quad \texttt{ioannis@automorph.io} \quad \textcopyright\ 2026 \quad CC-BY 4.0}
  \renewcommand{\headrulewidth}{0pt}
  \renewcommand{\footrulewidth}{0pt}
}

\title{Navier--Stokes as a Layer-Dissolving\\ Membrane Theorem}
\author{Ioannis Tsiokos\,\orcidlink{0009-0009-7659-5964}}
\date{\today}
\hypersetup{
  pdftitle={Navier--Stokes as a Layer-Dissolving Membrane Theorem},
  pdfauthor={Ioannis Tsiokos},
  hidelinks
}

\begin{document}

\maketitle
\thispagestyle{firstpage}

\begin{abstract}
We apply the adequacy-residual calculus of
\citet{Tsiokos2026XiCompanion} and the layer-dissolving membrane
endpoint of \citet{Tsiokos2026NeedleKiller} to the smooth-data
three-dimensional incompressible Navier--Stokes Cauchy problem.  The
argument fixes the Bradshaw--Grujic / Beale--Kato--Majda dissolving
readout \(\DBG\), measures the physical-native adequacy residual
\(\XiRes_C(\DBG\mid\Lphys)\), and propagates it along the doubling
Bradshaw--Grujic stages through the audited source-admission ledger
\(\Omsrc=\Omtr+\Ompr+\Omvisc+\Omtail+\Omstr\).  A no-go family rules
out packing, sparseness, marker-only, and metric-shadow shortcuts; a
Gevrey radius-window gate \(\lambdaJ\rho\gtrsim\log\lambdaJ\) isolates
the half-log endpoint obstruction; and a Littlewood--Paley transport
map \(\LPtr\colon L^2(P_j)\to H^{-s}\), \(s\ge q-\eps\), restores
summable predictive propagation.  Under the closure bundle
\(R_1,\ldots,R_9\) and the typed-context discipline, every
divergence-free initial datum \(u_0\in H^{\gamma_0}\) with
\(\gamma_0>5/2\) on the accepted smooth-data channel therefore
satisfies \(\Tmax(u_0)=\infty\).  The same closure records are then
recast as an AOR-instance membership witness for the smooth-data
carrier, using the audit-closure meta-theory of
\citet{Tsiokos2026AOR}, giving the record-level reading of ``the
carrier does not blow up in reality''; that recasting is presented
under the carrier-attached discharge-assignment discipline of the AOR
meta-theory; a data-bearing strengthening of the discharge-record
fields is outside the present scope.  We make no unconditional
Clay--Millennium claim, no generic critical-Besov data claim, no
criterion-native rewriting of \(\DBG\), and no claim of regularity
independent of the audit records.
\end{abstract}

\medskip
\noindent\textbf{Keywords.} Six Birds theory; Navier--Stokes regularity; layer-dissolving membrane; adequacy residual; Bradshaw--Grujic window; BKM continuation; Gevrey radius-window; Littlewood--Paley transport; smooth-data subcarrier; audit closure.



\section{Introduction}
\label{sec:intro}

We study the three-dimensional incompressible Navier--Stokes Cauchy
problem as a smooth-data regularity problem under an audit-closure
discipline.  The question is the classical one: whether a smooth
divergence-free datum can generate a finite-time loss of strong
regularity.  Under a structured closure bundle \(R_1,\ldots,R_9\) and
the typed-context discipline of \citet{Tsiokos2026XiCompanion} and
\citet{Tsiokos2026NeedleKiller}, every smooth divergence-free initial
velocity on the accepted smooth-data channel has infinite maximal
strong-solution lifespan.  The same records are also recast in
\Cref{sec:ns-aor-instance} as the AOR-membership witness for the
smooth-data carrier \(\mathsf C_{\mathrm{NS}}\) inside the
audit-closure meta-theory of \citet{Tsiokos2026AOR}, giving the
record-level reading that the carrier does not blow up in reality.
That recasting reframes the conclusion; it does not remove the
audit-record requirement.

Concretely, the Cauchy problem is posed for a velocity field \(u(x,t)\),
pressure \(p(x,t)\), and viscosity \(\nu>0\) on
\(\mathbb R^3\times[0,T)\):
\[
\begin{aligned}
  \partial_t u + (u\cdot\nabla)u + \nabla p &= \nu \Delta u,\\
  \nabla\cdot u &= 0,\\
  u(\cdot,0) &= u_0 .
\end{aligned}
\]
Throughout, smooth data means a divergence-free initial velocity
\(u_0\in H^{\gamma_0}\) with \(\gamma_0>5/2\); above this Sobolev
threshold Kato's classical theory~\citep{Kato1972} gives a unique
strong solution
\[
  u\in C\bigl([0,\Tmax(u_0)),H^{\gamma_0}\bigr)
\]
with maximal lifespan \(\Tmax(u_0)\in(0,\infty]\).  The regularity
question is whether \(\Tmax(u_0)=\infty\) for every such datum.

Two classical continuation criteria frame the question.  A standard
Beale--Kato--Majda-type continuation criterion for strong
Navier--Stokes flows~\citep{BealeKatoMajda1984,MajdaBertozzi2002,RobinsonRodrigoSadowski2016}
says that a strong solution continues past time \(T\) whenever the time
integral of the \(L^\infty\)-norm of its vorticity
\(\omega=\nabla\times u\) is finite,
\[
  \int_0^T \|\omega(\cdot,t)\|_{L^\infty}\,dt < \infty .
\]
The Bradshaw--Grujic frequency-window
criterion~\citep{BradshawGrujic2017} provides a complementary
frequency-localized route: continuation holds when a fixed observable,
read on a doubling sequence of Littlewood--Paley windows with
\(\lambda_{j+1}=2\lambda_j\), is uniformly bounded.  We therefore ask
whether the fixed frequency-window readout can be controlled along the
stages on which the strong solution is defined.

Our approach measures this question by a finite-dimensional
Schur-complement residual~\citep{Zhang2005SchurApplications}.  At a
Bradshaw--Grujic stage \(J\) and time \(t\), let \(C_{J,t}\) denote the
positive audit form on the smooth-data carrier, the energy form against
which observables are measured.  For a target observable \(D\) and a
native probe family \(L\), meaning a declared linear family of
measurements on the carrier, the adequacy residual is
\[
  \Xi_C(D\mid L)
  =
  K_{DD}-K_{DL}K_{LL}^{\dagger}K_{LD},
\]
where the \(K_{\ast\ast}\) are the Gram blocks determined by \(D\),
\(L\), and \(C\), and \(\dagger\) denotes the Moore--Penrose
pseudoinverse~\citep{Penrose1955GeneralizedInverse}.  The residual is a
positive operator on the response space measuring the part of \(D\)
unexplained by \(L\); it vanishes exactly when \(L\) captures \(D\)
after quotienting the zero-audit directions of \(C\).  The full
residual calculus is developed in \citet{Tsiokos2026XiCompanion}; we
use only this Schur-complement interface and its finite-dimensional
bookkeeping.

For Navier--Stokes the target observable is the fixed
Bradshaw--Grujic dissolving readout \(\DBG\): the same
frequency-localized quantity whose uniform boundedness on the doubling
Littlewood--Paley sequence \(\lambda_{j+1}=2\lambda_j\) gives the
classical continuation channel.  The native family is \(\Lphys\), a
finite linear family of Littlewood--Paley shell projections and Sobolev
(Bessel-potential) component probes on the Navier--Stokes carrier.  The
distinction matters: \(\Lphys\) is physically motivated and linear by
construction, not a relabeling of \(\DBG\).  The residual question is
whether \(\Xi_C(\DBG\mid\Lphys)\) can be bounded along the doubling
stages using audited Navier--Stokes records and a controlled list of
admitted source terms.

The principal regularity statement is the following.  For every
divergence-free datum \(u_0\in H^{\gamma_0}\) with \(\gamma_0>5/2\) in
the accepted smooth-data channel, the closure bundle
\(R_1,\ldots,R_9\) and the typed-context discipline imply
\(\Tmax(u_0)=\infty\).  The qualifier \emph{accepted smooth-data
channel} refers to the Sobolev data class admitted by the records used
in the proof; it is not an assertion about every weak or critical-data
Navier--Stokes carrier.

The proof is organized around the nine proof-carrying records collected
along the doubling Bradshaw--Grujic stages.  Each record is a named
bundle whose components hold the assumed PDE hypotheses and the
framework bookkeeping at the corresponding stage.  The records collect
the imported PDE results --- the Beale--Kato--Majda continuation
criterion, the Bradshaw--Grujic frequency-window doubling criterion,
Kato local well-posedness, half-log analyticity, the
Littlewood--Paley / Bony--Chemin--Danchin calculus, and the
Sobolev/Besov embedding chain --- together with the bookkeeping needed
to compose the residual argument: the source-admission ledger, the
audited transport map, the LP-transport stage operator, and the
typed-context deliverability bridge.  The typed-context discipline is
the requirement that the carrier, probes, transport maps, audit form,
and imported hypotheses remain declared throughout the comparison.

The proof instantiates the layer-dissolving membrane endpoint theorem
of \citet{Tsiokos2026NeedleKiller} on the Navier--Stokes carrier and
supplies the closure records as its hypothesis stack; the endpoint
converts that audited closure into the regularity conclusion.  The
audit-record requirement is essential at every stage; we make no
standalone Clay--Millennium
claim~\citep{Fefferman2000NavierStokes} independent of the closure
discipline and no data-bearing R-record claim beyond the present
record presentation.

\Cref{sec:ns-aor-instance} recasts the same closure records as the
AOR-membership witness for \(\mathsf C_{\mathrm{NS}}\) inside the
audit-closure meta-theory of \citet{Tsiokos2026AOR}.  Under that
recasting, the carrier belongs to the refinement-stable AOR class and
the regularity conclusion reads as the audit-closed statement that
the carrier does not blow up in reality.  This is a presentational
recasting: no new PDE input is introduced, the audit content is
unchanged, and the AOR-instance theorem
\Cref{thm:ns:in-refstableaor} is presented under the
carrier-attached discharge-assignment discipline of
\citet{Tsiokos2026AOR}, with a data-bearing strengthening of
the discharge-record fields stated as framework future work.

The paper is organized as follows.  \Cref{sec:xi-membrane} recalls the
adequacy-residual calculus from \citet{Tsiokos2026XiCompanion} and the
layer-dissolving membrane endpoint from
\citet{Tsiokos2026NeedleKiller}, and specializes both to the
Navier--Stokes carrier.  \Cref{sec:ns-carrier} fixes the carrier, the
fixed readout \(\DBG\), the physical-native family \(\Lphys\), the
BG-windowed exact package, and null-mode legality.  \Cref{sec:imports}
records the six imported PDE results and rules out the criterion-native
shortcut (NG1).  \Cref{sec:adequacy-cascade} develops the
source-admission ledger \(\Omsrc\) and the stage-to-stage recurrence
for the physical-native residual along the doubling stages.
\Cref{sec:no-go} states the remaining no-go family NG2--NG10,
excluding scalar-energy, packing, sparse-shadow, marker-only,
metric-shadow, moving-window, half-log endpoint, available-analyticity,
and raw-nonlinearity shortcuts.  \Cref{sec:gevrey-lp} develops the
Gevrey radius-window branch and the Littlewood--Paley transport rescue.
\Cref{sec:main-theorem} assembles \(R_1,\ldots,R_9\) and concludes
the framework-deliverable regularity theorem.
\Cref{sec:ns-aor-instance} recasts the closure records as the
AOR-membership witness for the smooth-data carrier.
\Cref{sec:status} closes with the external interpretation and the
locked nonclaim boundary.

 

\section{Adequacy Residual and Layer-Dissolving Membranes}
\label{sec:xi-membrane}

This section recalls the two framework tools used in the rest of the
paper: the adequacy-residual calculus of
\citet{Tsiokos2026XiCompanion} and the layer-dissolving membrane
endpoint of \citet{Tsiokos2026NeedleKiller}.

The residual calculus works with a finite-dimensional carrier equipped
with a positive audit form \(C\), a target observable \(D\), and a
native probe family \(L\) (a declared linear collection of measurements
on the carrier).  The adequacy residual is the Schur complement
\[
  \Xi_C(D\mid L)
  =
  K_{DD}-K_{DL}K_{LL}^{\dagger}K_{LD},
\]
where the \(K_{\ast\ast}\) are the Gram blocks of \(D\) and \(L\) under
\(C\) and \(\dagger\) is the Moore--Penrose pseudoinverse.

Geometrically, \(\Xi_C(D\mid L)\) is the squared distance from \(D\)
to the range of \(L\) in the audited Gram geometry; equivalently, it
is the part of \(D\) orthogonal to the best linear native explanation
under \(C\).  The optimal residual identity makes this precise: among
all linear native explanations of \(D\), the one obtained by projecting
onto the range of \(L\) gives the smallest residual in the Loewner
(positive-semidefinite) order~\citep{Bhatia2007PositiveDefinite}, and
that minimum is \(\Xi_C(D\mid L)\).  Exact adequacy is the zero-residual case:
\(\Xi_C(D\mid L)=0\) iff \(L\) explains \(D\) on the quotient by the
zero-audit directions of \(C\).

The calculus supplies four further rules used below.  Adequacy
promotion combines a native budget with a residual budget into a
dissolving budget, the kind of bound consumed by the endpoint theorem
recalled below.  The chain rule states that enlarging the native
family can only reduce the unexplained part:
\(\Xi_C(D\mid L')\preceqq\Xi_C(D\mid L)\) whenever
\(L\subseteq L'\).  Data processing records the behaviour under
declared maps: post-processing the target response transports the
residual by the response map; coarsening the native probes can only
leave the residual unchanged or larger.  Finally, the blind-spot
witness theorem identifies, whenever the residual violates an external
budget, a component of \(D\) invisible to \(L\) that accounts for the
violation.

The second imported tool is the membrane endpoint of
\citet{Tsiokos2026NeedleKiller}.  A predictive membrane is the
stagewise budget condition that the native predictive currency
\(\widehat K_j\) at every stage \(j\) fits inside a target dissolving
budget \(\Theta\),
\[
  \widehat K_j \preceqq \Theta ,
\]
i.e., the native explanations of the target observable stay inside one
declared budget across all stages.  A predictive needle is the
opposite phenomenon: a stagewise budget violation that would obstruct
the continuation conclusion.  The membrane is precisely the structured
condition ruling out such violations.

The layer-dissolving membrane endpoint converts this stagewise
condition into a final exclusion statement.  Under the all-six
discipline --- the record-keeping requirement that every comparison
declare its carrier, audit form, native probes, target observable,
transport maps, and imported records --- the existence of a predictive
membrane for the dissolving readout, together with audited transport
budgets and a typed-context deliverability bridge, implies the
no-predictive-needles conclusion.  We import this endpoint as a
previously proved result; the proof is in
\citet{Tsiokos2026NeedleKiller}.

\citet{Tsiokos2026NeedleKiller} also supplies a summable predictive
propagation lemma: when a stage budget contracts by a summable
envelope, the stagewise membrane condition propagates to a global
no-needle conclusion.  This ancillary import is used later when the
transported budget is made contractive stage by stage.

With the calculus and endpoint in hand, the Navier--Stokes
specialization is immediate.  The target observable is \(D=\DBG\), the
fixed dissolving readout, and the native family is \(L=\Lphys\), the
physical-native collection of Littlewood--Paley shell projections and
Sobolev component probes (formal definitions are given in
\Cref{sec:ns-carrier}).  The object to keep in view is the stagewise
residual
\[
  \Xi_{C_{J,t}}(\DBG\mid\Lphys)
\]
on the doubling stages \(\lambda_{j+1}=2\lambda_j\).  Bounding this
residual stage-by-stage is the central technical task of the paper.
If such a bound holds --- assembled below through the source-admission
ledger \(\Omsrc\), the no-go sieve, and the Gevrey radius-window with
Littlewood--Paley transport rescue --- then the layer-dissolving
endpoint applied to the resulting closure records \(R_1,\ldots,R_9\)
gives \(\Tmax(u_0)=\infty\) on the accepted smooth-data subcarrier
(\Cref{sec:main-theorem}).

 

\section{The Navier--Stokes Carrier and the BG/BKM Dissolving Readout}
\label{sec:ns-carrier}

Fix a divergence-free initial datum
\(u_0\in H^{\gamma_0}(\mathbb R^3)\) with \(\gamma_0>5/2\) and let
\[
  u\in C\bigl([0,\Tmax(u_0)),H^{\gamma_0}\bigr)
\]
be the unique strong solution given by Kato's local well-posedness
theorem.

The Bradshaw--Grujic indexing is parametrised by a time-dependent
stage \(J=J(t)\) tracking the active frequency scale: advancing the
stage doubles the scale,
\[
  \lambda_{j+1}=2\lambda_j ,
\]
and \(P_{J(t)}\) denotes the active Littlewood--Paley shell at time
\(t\)~\citep{Stein1970,Grafakos2014,BahouriCheminDanchin2011}, with
frequency scale \(\lambdaJ(t)\).  For each \(t\in[0,\Tmax(u_0))\) and
stage \(J\) we define the audited finite-dimensional carrier
\(E_{J,t}\) to be the span of the active shell components
\(\{P_j u(\cdot,t):j\le J\}\); it is a stage-indexed snapshot of the
strong solution.

The audit form \(C_{J,t}\) is the positive semidefinite form on
\(E_{J,t}\) against which observables are measured: in Fourier
variables, the \(H^{\gamma_0}\) Sobolev weight \((1+|\xi|^2)^{\gamma_0}\)
restricted to the modes retained by \(j\le J\).  The form may have a
non-trivial kernel \(\Ker C_{J,t}\) of zero-audit directions; residuals
are computed after quotienting by this kernel, and the legality
conditions below require the declared probes to vanish on it.

\begin{definition}[BG/BKM dissolving readout]\label{def:ns:bg-readout}
For each Bradshaw--Grujic stage \(J\) and time
\(t\in[0,\Tmax(u_0))\), the \emph{dissolving readout} is a
linear map
\[
  \DBG_{J,t}\colon E_{J,t}\to Z_J ,
\]
where \(Z_J\) is the finite-dimensional response space of
windowed Fourier values at frequency scale
\(\lambdaJ(t)\).  The readout evaluates the windowed
observable on \(P_{J(t)}\) appearing in the
Bradshaw--Grujic continuation criterion: typically an
\(L^p\) or \(L^\infty\) windowed norm of the vorticity or
velocity.  The right-support constraint
\(\DBG_{J,t}=\DBG_{J,t}\cdot P_{J(t)}\) records that
\(\DBG_{J,t}\) sees only the active shell at stage \(J\).
\end{definition}

The family \(\{\DBG_{J,t}\}\) is fixed throughout the paper and is
\emph{not} replaced by a criterion-native surrogate.  A criterion-native
surrogate is any relabeling such as \(\LBG:=\DBG\) that declares the
native probe to be the target itself; such a relabeling forces
\(\Xi_C(\DBG\mid\LBG)=0\) by definition without explaining the
Bradshaw--Grujic observable through independent physical measurements.
\Cref{sec:imports} records this as a no-go (NG1).

The same readout serves both classical continuation criteria from
\Cref{sec:intro}: the Bradshaw--Grujic frequency-window criterion
consumes its uniform boundedness on doubling stages, while the
Beale--Kato--Majda criterion consumes the implied vorticity-time
integrability.

\begin{definition}[Physical-native probe family]\label{def:ns:physical-native-probes}
For each Bradshaw--Grujic stage \(J\) and time \(t\), the
\emph{physical-native probe family} is the finite linear map
\[
  \Lphys_{J,t}\colon E_{J,t}\to Y ,
\]
where \(Y\) collects the values of the declared physical
measurements.  Each component of \(\Lphys_{J,t}\) is either
a Littlewood--Paley shell projection \(P_j u(\cdot,t)\)
with \(j\le J\) in an \(H^{\gamma_0}\)-compatible norm, or
a classical Sobolev (Bessel-potential) component probe such
as an \(H^s\)-component of the velocity with
\(s\le\gamma_0\).  The native Schur block
\(\Lphys_{J,t}C_{J,t}^{\dagger}(\Lphys_{J,t})^*\) recovers
the kinetic-energy and enstrophy budgets visible at the
Bradshaw--Grujic window resolution; energy and enstrophy
themselves are quadratic in the carrier data and are
therefore not native probes (a no-go recorded in
\Cref{sec:no-go}).
\end{definition}

The residual studied below is \(\Xi_{C_{J,t}}(\DBG\mid\Lphys)\).

\begin{definition}[BG-windowed exact package]\label{def:ns:bg-package}
At each Bradshaw--Grujic stage \(j\), the \emph{BG-windowed
exact package} \(\Gamma_j\) collects the admissible
BG-window probes (the relevant components of
\(\Lphys_{J,t}\) and the windowed value of \(\DBG_{J,t}\)
on the active shell \(P_j\)) together with the structural
constraints they must preserve:
\begin{enumerate}[label=(\alph*)]
  \item the divergence-free condition \(\nabla\cdot u=0\);
  \item the Helmholtz-projected carrier
    \citep{Galdi2011} on which the probes act;
  \item the boundary arrangement relating \(P_j\) to
    \(P_{j+1}\);
  \item the doubling relation
    \(\lambda_{j+1}=2\lambda_j\).
\end{enumerate}
The package is \emph{exact} in the sense of
\citet{Tsiokos2026NeedleKiller}: it transports the audited
sub-problem from one BG window to the next without
altering the energy, response, or residual information.
\end{definition}

\begin{definition}[Null-mode legality of the NS carrier]\label{def:ns:null-mode-legal-ns}
The \emph{NS carrier} \((E_{J,t},C_{J,t},\Lphys,\DBG)\) is
\emph{null-mode legal} when the physical-native probes and
the dissolving readout both vanish on the audit kernel:
\[
  \Lphys_{J,t}\big|_{\Ker C_{J,t}}=0,
  \qquad
  \DBG_{J,t}\big|_{\Ker C_{J,t}}=0 .
\]
Equivalently, both families descend to the legal energy
quotient \(E_{J,t}/\Ker C_{J,t}\), on which the audit form
is positive definite and the adequacy residual
\(\Xi_{C_{J,t}}(\DBG\mid\Lphys)\) of
\citet{Tsiokos2026XiCompanion} is well-defined.  Null-mode
legality is built into the carrier by construction; it is
not a regularity assumption to be verified after the
residual has been formed.
\end{definition}

 

\section{Literature Imports and the Criterion-Native Gap}
\label{sec:imports}

This section records six PDE imports --- the Beale--Kato--Majda
continuation criterion, the Bradshaw--Grujic frequency-window criterion,
Kato local well-posedness, half-log analyticity, the
Littlewood--Paley / Bony--Chemin--Danchin calculus, and the
Sobolev/Besov embedding chain --- and the criterion-native gap that
motivates the residual cascade.  Each import enters as a proof-carrying
record: a structured assumption whose stated hypotheses expose the
classical conclusion, with the conclusion supplied by the PDE
literature rather than derived here.

\begin{definition}[Beale--Kato--Majda continuation criterion]\label{def:ns:bkm-import}
The \emph{Beale--Kato--Majda continuation criterion}
\citep{BealeKatoMajda1984,MajdaBertozzi2002,RobinsonRodrigoSadowski2016}
is the imported record asserting that for a strong solution
\(u\) of the Navier--Stokes equations on \([0,T)\) with
\(T\le\Tmax(u_0)\) and \(\omega=\nabla\times u\),
\[
  \int_0^T \|\omega(\cdot,t)\|_{L^\infty}\,dt < \infty
  \quad\Longrightarrow\quad
  \Tmax(u_0)>T .
\]
The record is consumed at its declared status, with the
continuation conclusion supplied by the cited references
rather than derived here.
\end{definition}

\begin{definition}[Bradshaw--Grujic frequency-window continuation criterion]\label{def:ns:bg-import}
The \emph{Bradshaw--Grujic frequency-window continuation
criterion} \citep{BradshawGrujic2017} is the imported
record asserting that, for the dissolving readout of
\Cref{def:ns:bg-readout} and doubling stages
\(\lambda_{j+1}=2\lambda_j\),
\[
  \sup_j\sup_{t\in I_j}\|\DBG_{j,t}u\|_{Z_j}<\infty
  \quad\Longrightarrow\quad
  \Tmax(u_0)=\infty .
\]
The record is consumed at its declared status; this paper
constructs an audited stage-by-stage bound on
\(\Xi_{C_{J,t}}(\DBG\mid\Lphys)\) that delivers this
hypothesis.
\end{definition}

\begin{definition}[Kato local well-posedness on the smooth-data subcarrier]\label{def:ns:kato-import}
The \emph{Kato local well-posedness import}
\citep{Kato1972,FujitaKato1964,RobinsonRodrigoSadowski2016}
is the imported record asserting that for divergence-free
smooth Sobolev data on \(\mathbb R^3\) with
\(\gamma_0>5/2\),
\[
  \nabla\cdot u_0=0,\quad u_0\in H^{\gamma_0}(\mathbb R^3)
  \quad\Longrightarrow\quad
  \exists\,\Tmax(u_0)\in(0,\infty]\text{ and a unique }
  u\in C\bigl([0,\Tmax(u_0)),H^{\gamma_0}\bigr).
\]
The carrier objects \(E_{J,t}\), \(C_{J,t}\),
\(\Lphys_{J,t}\), and \(\DBG_{J,t}\) of
\Cref{sec:ns-carrier} are all attached to this
strong-solution arc.
\end{definition}

\begin{definition}[Half-log analyticity import]\label{def:ns:wang-halflog-import}
The \emph{half-log analyticity import}, in the
Foias--Temam tradition \citep{FoiasTemam1989} refined to
half-log form by \citet{HerbstSkibsted2011} and by the
preprint \citep{Wang2025AnalyticityRefined}, is the
imported record asserting that restart data in the joint
space
\[
  H^{\gamma_0}\cap\dot B^{-\eps}_{\infty,\infty},
  \qquad \gamma_0>3/2-\eps,
\]
yields a lower bound on the spatial analyticity radius
\(\rho(t)\): on BG-compatible stages,
\[
  \rho(t)\gtrsim \frac{\log\lambdaJ(t)}{\lambdaJ(t)} .
\]
The record is consumed at its declared status.
\end{definition}

\begin{definition}[Littlewood--Paley / Bony--Chemin--Danchin calculus]\label{def:ns:lp-bcd-import}\leavevmode\par\noindent
The \emph{Littlewood--Paley / Bony--Chemin--Danchin
paradifferential calculus import}
\citep{Bony1981,BahouriCheminDanchin2011} on
\(\mathbb R^3\) is the imported record supplying the
shell-orthogonality identities, commutator estimates, and
operator-norm bounds between Littlewood--Paley shells
\(P_j\) and smooth multipliers used downstream.  The
operator-norm decay
\(\|\LPtr\|^2\simeq 2^{-2sj}\) of the LP-transport stage
operator (\Cref{def:ns:lp-transport}) is recovered from
this calculus.
\end{definition}

\begin{definition}[Sobolev/Besov embedding chain]\label{def:ns:sobolev-besov-embedding-import}
The \emph{Sobolev/Besov embedding chain import}
\citep{Triebel1983,BahouriCheminDanchin2011} is the
imported record asserting that the smooth Sobolev restart
data lifts to the Besov component required by the half-log
import:
\[
  H^{\gamma_0}\hookrightarrow \dot B^{-\eps}_{\infty,\infty},
  \qquad \gamma_0>3/2,\quad \eps>0\text{ admissible}.
\]
Together with \Cref{def:ns:lp-bcd-import} and
\Cref{def:ns:wang-halflog-import}, this import forms the
analytic infrastructure consumed by the Gevrey
radius-window branch and the LP-transport rescue of
\Cref{sec:gevrey-lp}.
\end{definition}

We close this section with an algebraic equality and the no-go it
motivates.

\begin{theorem}[Criterion-native equality]\label{thm:ns:criterion-native-equality}
Set \(\LBG:=\DBG\).  Then on the legal quotient of \(C_{J,t}\),
\[
  \Xi_{C_{J,t}}(\DBG\mid \LBG)=0 .
\]
\end{theorem}

\begin{proof}
Substituting \(L=\LBG=\DBG\) into the Schur-complement definition gives
\(K_{LD}=K_{DD}\) and
\(K_{DL}K_{LL}^{\dagger}K_{LD}=K_{DD}K_{DD}^{\dagger}K_{DD}=K_{DD}\)
on the legal quotient of \(C_{J,t}\).  Hence
\(\Xi_{C_{J,t}}(\DBG\mid\LBG)=K_{DD}-K_{DD}=0\) on that legal
quotient.\footnote{Verified in Lean as \texttt{SixBirdsNeedles.NS.CriterionNativeGap.criterionNativeEquality}; see \Cref{subsec:paper3-per-label-coverage}.}
\end{proof}

\begin{theorem}[No-go NG1: criterion-native circularity]\label{thm:ns:no-go-criterion-native}
The criterion-native equality obtained from \(\LBG:=\DBG\) does not
bound the physical-native residual.  There are audited finite data for
which
\[
  \Xi_{C_{J,t}}(\DBG\mid\LBG)=0
  \quad\text{but}\quad
  \Xi_{C_{J,t}}(\DBG\mid\Lphys)\ne 0 .
\]
Thus criterion-native labeling is not a witness for physical-native
adequacy.
\end{theorem}

\begin{proof}
A simple finite witness suffices.  Take the audit form \(C=I_2\), the
target row \(\DBG=(1,0)\), and the physical-native row
\(\Lphys=(0,1)\): the criterion-native residual vanishes after the
relabeling \(\LBG:=\DBG\), while the residual against \(\Lphys\)
remains nonzero.\footnote{Mechanized in Lean as
\texttt{SixBirdsNeedles.NS.CriterionNativeGap.noGoCriterionNative}
as a partial mechanization: the Lean theorem proves the
specific 2-dimensional finite witness
(\(C=I_2\), \(\DBG=(1,0)\), \(\Lphys=(0,1)\)) using identity
pseudoinverses; the general statement is proved here.  See
\Cref{subsec:paper3-per-label-coverage}.}
\end{proof}

The point is elementary: relabeling the target as native does not
extract any information from independent physical measurements.  The
remaining sections build an audited stage-by-stage bound on the
physical-native residual, since the criterion-native shortcut is
ruled out.

 

\section{The Physical-Native Adequacy Cascade}
\label{sec:adequacy-cascade}

This section builds the stage-by-stage recurrence that propagates the
physical-native residual \(\Xi_{C_{J,t}}(\DBG\mid\Lphys)\) along the
doubling Bradshaw--Grujic stages.  The recurrence is audited by a
positive operator-valued ledger \(\Omsrc\), which records the new
contributions entering the dissolving response space at each stage.

\begin{definition}[Source-admission ledger]\label{def:ns:source-admission-ledger}
The \emph{source-admission ledger} is the positive
operator-valued sum
\[
  \Omsrc = \Omtr + \Ompr + \Omvisc + \Omtail + \Omstr ,
\]
with five positive semidefinite components
\[
  \Omtr\succeq 0,\qquad
  \Ompr\succeq 0,\qquad
  \Omvisc\succeq 0,\qquad
  \Omtail\succeq 0,\qquad
  \Omstr\succeq 0,
\]
hence \(\Omsrc\succeq 0\).  The summands record the
contributions entering the dissolving response space at each
stage from, respectively:
\begin{enumerate}[label=(\alph*)]
  \item the nonlinear transport \((u\cdot\nabla)u\) as a Schur
    block on the active shells (\(\Omtr\));
  \item the pressure \(\nabla p\) after the Helmholtz/Leray
    projection (\(\Ompr\));
  \item viscosity \(\nu\Delta u\) as a damping block on each
    shell (\(\Omvisc\));
  \item the high-frequency tail capturing the interaction of
    visible shells with discarded shells \(j>J\)
    (\(\Omtail\));
  \item vortex stretching \((\omega\cdot\nabla)u\) tested
    against the physical-native probes (\(\Omstr\)).
\end{enumerate}
Each summand is owned by an imported PDE record from
\Cref{sec:imports}: \(\Omtr\) and \(\Ompr\) by the
Littlewood--Paley / Bony--Chemin--Danchin commutator and
Helmholtz-projection estimates
(\Cref{def:ns:lp-bcd-import}); \(\Omvisc\) by the viscous
damping at scale \(\lambdaJ(t)\) from \(\nu\Delta\);
\(\Omtail\) by the Sobolev/Besov embedding combined with
the BG-window cutoff (\Cref{def:ns:sobolev-besov-embedding-import});
and \(\Omstr\) by the vortex-stretching estimate against
the physical-native probes.
\end{definition}

For a time step \(h>0\) with \(t+h<\Tmax(u_0)\), let
\(\mathcal T_h:E_{J,t}\to E_{J,t+h}\) denote the audited transport map
between the stage-\(J\) carriers at \(t\) and \(t+h\): it is the
linear semigroup-style propagation of probe data across the step, in
the linear-source-corrected evolution of the active shells.  By
declaration \(\mathcal T_h\) descends to the legal energy quotient,
\[
  \mathcal T_h(\Ker C_{J,t})\subseteq \Ker C_{J,t+h},
\]
so the induced map
\(\bar{\mathcal T}_h:E_{J,t}/\Ker C_{J,t}\to E_{J,t+h}/\Ker C_{J,t+h}\)
is well-defined and the Loewner recurrence below is typed on the
legal quotient throughout; on the smooth-data Sobolev carrier
\(C_{J,t}\) is positive-definite on \(E_{J,t}\), so
\(\Ker C_{J,t}=\{0\}\) and the descent is automatic.  The carry-forward
constant \(C_{\rm old}\ge 0\) records how much of the old residual
survives this transport; it is determined by the Littlewood--Paley /
Bony--Chemin--Danchin calculus and the Sobolev/Besov embedding
records, not a free parameter.  These conventions are inherited from
the residual calculus of \citet{Tsiokos2026XiCompanion}.

\begin{theorem}[Source-admission recurrence]\label{thm:ns:source-admission-recurrence}
Under the null-mode legality and source-admission acceptance records,
for every BG stage \(J\), time \(t\), and time step \(h>0\) with
\(t+h<\Tmax(u_0)\),
\[
  \Xi_{C_{J,t+h}}\bigl(\DBG_{J,t+h}\mid \Lphys_{J,t+h}\bigr)
  \preceqq
  C_{\rm old}\,
  \mathcal T_h\,
  \Xi_{C_{J,t}}\bigl(\DBG_{J,t}\mid \Lphys_{J,t}\bigr)\,
  \mathcal T_h^*
  + \Omsrc .
\]
\end{theorem}

\begin{proof}
The residual at time \(t+h\) is bounded above by transporting the
residual at time \(t\) through \(\mathcal T_h\) (paying the
carry-forward constant \(C_{\rm old}\)) and adding the newly admitted
sources, which are the five positive components of \(\Omsrc\) listed
above.\footnote{Verified in Lean as \texttt{SixBirdsNeedles.NS.Cascade.sourceAdmissionRecurrence}; see \Cref{subsec:paper3-per-label-coverage}.}
\end{proof}

The recurrence organises downstream into a residual tree.  Successive
refinement steps either pay a piece using an imported record, expose
a piece as a shortcut excluded by the no-go family, or pass the
remaining obligation to a later analytic mechanism:
\[
  \Xi_C(\DBG\mid\Lphys)
  \longrightarrow
  \Omsrc
  \longrightarrow
  \Omega_{\rm gap}
  \longrightarrow
  \Omega_{\rm rem}
  \longrightarrow
  \Omega_{\rm bb}
  \longrightarrow
  \Omega_{\rm mm}
  \longrightarrow
  \Omamp
  \longrightarrow
  \Omend .
\]
The ledger \(\Omega_{\rm gap}\) records the gap between the
transported old residual and the controlled native budget after
Beale--Kato--Majda, Bradshaw--Grujic, and Kato are applied;
\(\Omega_{\rm rem}\) is the remainder after the half-log analyticity,
Littlewood--Paley / Bony--Chemin--Danchin, and Sobolev/Besov records
have been folded in.  The remaining ledgers identify rather than pay
obstructions: \(\Omega_{\rm bb}\) collects the part of the residual
localised on geometrically or frequency-wise bad regions;
\(\Omega_{\rm mm}\) collects contributions whose sector or probe type
does not match \(\Lphys\); \(\Omamp\) is the amplitude obligation left
after the public-shadow shortcuts have been excluded; and \(\Omend\)
is the endpoint obstruction at the close of the cascade.  The no-go
family of \Cref{sec:no-go} excludes proposed shortcut payments ---
scalar energy, packing, sparse-shadow data, marker payloads, metric
shadows, moving windows, available half-log analyticity alone, raw
nonlinear observables --- but does not pay these ledgers; the
remaining amplitude and endpoint obligations are discharged by the
Gevrey radius-window branch and Littlewood--Paley transport rescue of
\Cref{sec:gevrey-lp}.

 

\section{Framework-Derived No-Go Theorems}
\label{sec:no-go}

This section records the no-go family constraining attempts to pay
the residual ledgers of \Cref{sec:adequacy-cascade}.  Each no-go
excludes a category of shortcut payment that uses data of the wrong
type.  The no-gos do not pay any ledger; they restrict admissible
payments so that the Gevrey radius-window and Littlewood--Paley
transport rescue of \Cref{sec:gevrey-lp} carries the remaining
analytic obligation.  The criterion-native circularity result NG1 was
recorded in \Cref{sec:imports}; here we record NG2--NG10.

\begin{theorem}[No-go NG2: scalar energy / enstrophy category error]
\label{thm:ns:no-go-scalar-energy}
Quadratic scalar budgets such as kinetic energy \(\|u\|_{L^2}^2\),
enstrophy \(\|\nabla u\|_{L^2}^2\), or any positive scalar functional
that is quadratic in the carrier variable, are not lawful native Schur
probes.  The adequacy residual
\[
  \Xi_C(D\mid L)=K_{DD}-K_{DL}K_{LL}^{\dagger}K_{LD}
\]
is defined for a linear native probe family \(L:E\to Y\).  A quadratic
scalar functional is not such a family, so
\(\Xi_C(\DBG\mid\text{scalar quadratic})\) is not a well-formed object
of the Schur-complement calculus until a lawful linear native family has
been declared and audited.
\end{theorem}

\begin{proof}
The point is a type-level category mismatch.  The Schur
blocks~\citep{HornJohnson2013MatrixAnalysis} \(K_{LL}\), \(K_{DL}\),
\(K_{LD}\) come from composing the audit form with linear maps out of
the carrier; a quadratic scalar does not provide such maps and cannot
be inserted as \(L\).  Energy and
enstrophy still enter the paper, but as derived budgets recovered from
the native currency block \(\Lphys C_{J,t}^{\dagger}(\Lphys)^*\) of
\Cref{sec:ns-carrier}: NG2 says the linear family \(\Lphys\), not the
quadratic scalars themselves, is the native probe family.\footnote{Verified in Lean as \texttt{SixBirdsNeedles.NS.NoGo.noGoScalarEnergy}; see \Cref{subsec:paper3-per-label-coverage}.}
\end{proof}

The next four no-gos share a common form.  Each rejects a
public-shadow datum: information visible in a geometric, combinatorial,
or weighted description of the flow but not in itself dominating the
fixed Bradshaw--Grujic readout in the Schur calculus.

\begin{theorem}[No-go NG3: packing / amplitude]
\label{thm:ns:no-go-packing-amplitude}
For a bad set with size scaling \(\sim r^{3-d}\) and amplitude scaling
\(\sim r^{-\alpha}\), the contribution to the BG readout scales as
\[
  r^{(3-d)/2-\alpha}.
\]
Packing alone suppresses this amplitude only when
\(\alpha < (3-d)/2\).  Outside that range, small-set sparseness does
not control \(\DBG\)-amplitude, and the residual requires either a
sector-matched amplitude bridge or a paid tail-smallness record.
\end{theorem}

\begin{proof}
A set may occupy little volume, but if its amplitude grows faster
than the square-root packing gain, the contribution seen by the fixed
BG window remains large.  Packing data therefore enter the cascade
only after sector matching, or after the missing tail-smallness is
paid in the source-admission ledger.\footnote{Verified in Lean as \texttt{SixBirdsNeedles.NS.NoGo.noGoPackingAmplitude}; see \Cref{subsec:paper3-per-label-coverage}.}
\end{proof}

\begin{theorem}[No-go NG4: sparse-shadow]
\label{thm:ns:no-go-sparse-shadow}
Imported sparse or intense sets do not pay the bad-cylinder /
bad-geometry ledger \(\Omega_{\rm bb}\) unless either the sparse field
is sector-contained in the bad-bad / mismatch sector, or a paid mismatch
record \(\Omega_{\rm mm}\succeq 0\) is supplied.
\end{theorem}

\begin{proof}
Sparseness tells us where a defect does not live but not which sector
contains the obstruction that the physical-native residual must
control.  The lawful options are explicit: sparse-shadow data is
admitted only when contained in the bad-bad or mismatch sector, or
when accompanied by a positive mismatch record.\footnote{Verified in Lean as \texttt{SixBirdsNeedles.NS.NoGo.noGoSparseShadow}; see \Cref{subsec:paper3-per-label-coverage}.}
\end{proof}

\begin{theorem}[No-go NG5: marker-only]
\label{thm:ns:no-go-marker-only}
Triad markers, bad-overlap masks, and similar marker-only fields identify
where a defect lies and what class it belongs to, but they do not
dominate the BG amplitude without a quantitative bridge.  Promotion from
marker payload to the amplitude ledger \(\Omamp\) requires either a new
native probe family admitted through
\Cref{thm:ns:source-admission-recurrence}, or a record-grade
\(M^{\rm bb}\) bound.
\end{theorem}

\begin{proof}
A marker is a locator, not a bound: it identifies a triad, an overlap
class, or a bad region without measuring the size of the component
seen by \(\DBG\).  The no-go therefore forces marker data to be coupled
to a genuine amplitude bridge --- either new linear probes admitted
through the source-admission recurrence, or a record-grade
\(M^{\rm bb}\) estimate that supplies the missing size information.\footnote{Verified in Lean as \texttt{SixBirdsNeedles.NS.NoGo.noGoMarkerOnly}; see \Cref{subsec:paper3-per-label-coverage}.}
\end{proof}

\begin{theorem}[No-go NG6: metric-shadow]
\label{thm:ns:no-go-metric-shadow}
Weighted or Gevrey-flavored metrics that shrink the apparent residual
do not control the fixed BG readout unless the normalized metric still
dominates the BG/BKM amplitude.  Geometry inputs from
Caffarelli--Kohn--Nirenberg partial
regularity~\citep{CaffarelliKohnNirenberg1982} or
Constantin--Fefferman vorticity-direction
geometry~\citep{ConstantinFefferman1993} act as metric shadows unless
a sector-matched amplitude bridge, an explicit dominance constant,
and a time gate are supplied as separate records.
\end{theorem}

\begin{proof}
The issue is normalisation.  A geometric metric may make the residual
look small by adopting weights adapted to the geometry, but the
continuation channel reads the fixed BG/BKM amplitude on the carrier
scale, not on the auxiliary metric.  To be admissible, the geometric
input must dominate the BG amplitude on the normalised carrier scale
and must state the time interval on which that dominance holds.\footnote{Verified in Lean as \texttt{SixBirdsNeedles.NS.NoGo.noGoMetricShadow}; see \Cref{subsec:paper3-per-label-coverage}.}
\end{proof}

\begin{theorem}[No-go NG7: moving-window / derivative-stack support-only]
\label{thm:ns:no-go-moving-window}
Finite Littlewood--Paley windows or derivative-stack windows, meaning
families such as \(\nabla^m P_Ju(\cdot,t)\) tracked for several
derivative orders \(m\), do not promote to a completed carrier closure
merely because their support is known to lie on moving sets.  Promotion
requires one of three additional records: a finite-termination record
showing that the moving support stabilizes after finitely many stages; a
fixed-ledger squeeze showing that the moving support is bounded by a
stage-independent ledger; or an exhaustive-ledger squeeze showing that
the moving support is contained in the union of audited ledgers.
\end{theorem}

\begin{proof}
Moving-support data records where an obstruction lives, not how its
budget is controlled.  A window may track the active frequency content
and a derivative stack may track several differentiated shells, but
support alone does not say that the obstruction terminates, fits a
fixed ledger, or is exhausted by audited ledgers.  The no-go forces a
termination or squeeze record on top of the support information.\footnote{Verified in Lean as \texttt{SixBirdsNeedles.NS.NoGo.noGoMovingWindow}; see \Cref{subsec:paper3-per-label-coverage}.}
\end{proof}

\begin{theorem}[No-go NG8: half-log endpoint obstruction]
\label{thm:ns:no-go-half-log-endpoint}
On the available half-log analyticity lane alone, the staged residual
budget has the form
\[
  B_j
  =
  C_B\,2^{2pj}\bigl(1+\mu\sqrt{L_j}\bigr)
  e^{-2c\mu\sqrt{L_j}},
  \qquad
  p=q-\eps,\qquad
  L_j=\log\lambdaJ .
\]
For positive constants \(C_B,\mu,c\) and \(p>0\), the stage ratio
satisfies
\[
  \frac{B_{j+1}}{B_j}
  =
  2^{2p}
  \frac{\bigl(1+\mu\sqrt{L_{j+1}}\bigr)
        e^{-2c\mu\sqrt{L_{j+1}}}}
       {\bigl(1+\mu\sqrt{L_j}\bigr)
        e^{-2c\mu\sqrt{L_j}}}
  \longrightarrow 2^{2p}>1 .
\]
Thus the half-log Gevrey budget under Bradshaw--Grujic
frequency-doubling fails the summable predictive-propagation hypothesis
unless an additional decay record is supplied.
\end{theorem}

\begin{proof}
The doubling relation gives
\(L_j=\log\lambdaJ=j\log 2+O(1)\), so
\(\sqrt{L_{j+1}}-\sqrt{L_j}=O(L_j^{-1/2})\to 0\).  The exponential
factor in the displayed ratio has exponent
\(-2c\mu(\sqrt{L_{j+1}}-\sqrt{L_j})\) and therefore tends to \(1\); the
polynomial factor \((1+\mu\sqrt{L_{j+1}})/(1+\mu\sqrt{L_j})\) also
tends to \(1\).  The only surviving multiplier is \(2^{2p}\), and
since \(p=q-\eps>0\) the limit exceeds \(1\): the staged budget grows
geometrically rather than forming a summable tail.\footnote{Verified in Lean as \texttt{SixBirdsNeedles.NS.NoGo.noGoHalfLogEndpoint}; see \Cref{subsec:paper3-per-label-coverage}.}
\end{proof}

\noindent
The obstruction is exactly what forces an additional decay mechanism
--- Littlewood--Paley transport, a strictly stronger full-log radius,
or a new cancellation.

\begin{theorem}[No-go NG9: available analyticity ceiling for the direct route]
\label{thm:ns:no-go-published-analyticity}
The available half-log analyticity lane meets the basic radius-window
gate \(\lambdaJ\rho(t)\gtrsim\log\lambdaJ\), but it falls short of the
direct-route predictive-contraction threshold by a factor of
\(\sqrt{\log\lambdaJ}\).  Concretely, NG8 shows that the staged budget
on the half-log lane has stage ratio tending to \(2^{2p}>1\).  Closing
the direct route therefore requires either the Littlewood--Paley
transport rescue of \Cref{thm:ns:lp-transport-rescue}, the proof of a
strictly stronger full-log analyticity radius, or an alternative
new-radius theorem outside the present scope.
\end{theorem}

\begin{proof}
The half-log import supplies the basic radius window, but NG8 shows
that this window does not by itself produce summable predictive
propagation on the direct route: contraction would require an
analyticity-radius improvement equivalent to a
\(\sqrt{\log\lambdaJ}\) factor, or else a different mechanism
contributing the missing decay.  In this paper that mechanism is the
Littlewood--Paley transport rescue of \Cref{sec:gevrey-lp}.\footnote{Verified in Lean as \texttt{SixBirdsNeedles.NS.NoGo.noGoPublishedAnalyticity}; see \Cref{subsec:paper3-per-label-coverage}.}
\end{proof}

\begin{theorem}[No-go NG10: raw nonlinearity illegal as Schur probe]
\label{thm:ns:no-go-raw-nonlinearity}
Raw nonlinear observables, including flux quantities, pressure
interactions, vortex stretching, and triad terms, cannot enter the
native family \(\Lphys\) as Schur-complement probes.  Lawful entry
requires one of the following:
\begin{enumerate}[label=(\roman*)]
\item tangent linearization around the audited solution arc, producing a
linear probe that may depend on the solution but is linear in the
perturbation;
\item interaction lifting to a higher-order linearization, for example
by expressing a triad as a linear probe on an augmented carrier that
explicitly resolves the interaction;
\item recording the nonlinear quantity as a source-admission entry in
\(\Omsrc\), so it is treated as part of the ledger rather than as a
native probe.
\end{enumerate}
\end{theorem}

\begin{proof}
This is the nonlinear analogue of the type discipline of
\Cref{thm:ns:no-go-scalar-energy}: the Schur complement is defined
for linear maps out of the carrier, and raw nonlinear quantities do
not provide the required Gram blocks.  Nonlinear physical content is
not excluded; it must enter through one of the three lawful interfaces
above.\footnote{Verified in Lean as \texttt{SixBirdsNeedles.NS.NoGo.noGoRawNonlinearity}; see \Cref{subsec:paper3-per-label-coverage}.}
\end{proof}

Taken together, NG1--NG10 exclude criterion-native circularity,
scalar quadratic budgets as native Schur probes, packing, sparseness,
markers, and metric shadows as sole bridges to amplitude,
moving-window or derivative-stack support without termination or
squeeze records, the available half-log lane as a summable
predictive-propagation mechanism on its own, the
\(\sqrt{\log\lambdaJ}\) shortfall for the direct contraction route,
and raw nonlinear observables as native Schur probes.  What remains
admissible is the typed data already declared: the physical-native
probes \(\Lphys\), the source-admission ledger, the transport
recurrence of \Cref{thm:ns:source-admission-recurrence}, the
Beale--Kato--Majda and Bradshaw--Grujic continuation records, Kato
local well-posedness, and the analytic imports of half-log
analyticity, Littlewood--Paley / Bony--Chemin--Danchin calculus, and
Sobolev/Besov embedding.  The remaining obstruction after the sieve is
\(\Omamp\), which is paid by the Gevrey radius-window branch and the
Littlewood--Paley transport rescue of \Cref{sec:gevrey-lp}.

 

\section{Gevrey Radius-Window Branch and LP-Transport Rescue}
\label{sec:gevrey-lp}

This section pays the remaining amplitude obligation \(\Omamp\) left by
the no-go sieve.  Working in the residual-calculus framework of
\citet{Tsiokos2026XiCompanion}, we introduce a lawful family of
Gevrey-normalised derivative probes, optimise the resulting dominance
constant over the derivative order, and express the result in terms of
the radius-window quantity \(\sigmaJ\).  We then bridge the half-log
endpoint obstruction through a Littlewood--Paley transport rescue on
the smooth-data subcarrier.

\begin{definition}[Gevrey-normalized derivative probes]
\label{def:ns:gevrey-probes}
Let \(\rho(t)\) be the spatial analyticity radius at time \(t\), as
supplied by the half-log analyticity import, and let \(\CG>0\) be the
Gevrey normalization constant.  For each derivative order \(m\ge 0\),
BG stage \(J\), and time \(t\), the Gevrey-normalized derivative probe
is
\[
  \widetilde M_{m,J,t}
  =
  \frac{\rho(t)^m}{\CG^m\,m!}\,\nabla^m P_{J(t)} .
\]
The family \(\{\widetilde M_{m,J,t}\}_{m\ge 0}\) is a lawful native
probe family on the BG-windowed carrier under the half-log analyticity
record.
\end{definition}

The weight \(\rho(t)^m/(\CG^m m!)\) is the standard Gevrey
normalisation: derivatives of order \(m\) are reweighted so that the
formal derivative series converges when the solution is analytic in a
strip of width \(\rho(t)\).  The dimensionless quantity attached to
the active BG window,
\[
  \sigmaJ(t)=\frac{\lambdaJ(t)\rho(t)}{\CG},
\]
measures the analyticity radius in units of the inverse active
frequency scale; it is the natural variable in which to state the
dominance estimate and radius-window comparison below.

\begin{definition}[LP transport stage operator]
\label{def:ns:lp-transport}
For each BG stage \(j\), the LP-transport stage operator is the linear
map
\[
  \LPtr\colon L^2(P_j)\longrightarrow H^{-s},
  \qquad s\ge q-\eps .
\]
On the smooth-data subcarrier it satisfies the operator-norm bound
\[
  \|\LPtr\|^2\simeq 2^{-2sj}.
\]
Here \(P_j\) is the Littlewood--Paley shell at frequency scale \(2^j\),
and \(s\) is the response exponent.
\end{definition}

The decay \(2^{-2sj}\) is the operator-norm content of the
Littlewood--Paley / Bony--Chemin--Danchin calculus specialised to the
BG window, consumed here as an audited input.  The map \(\LPtr\) takes
a budget on the \(j\)-th shell and transports it into a response
budget --- the size of the obstruction measured in \(H^{-s}\).  The
negative Sobolev codomain is natural because the transported object is
a response, not a function on the carrier.  Its
\(2^{-2sj}\) decay, combined with the Gevrey dominance and analyticity
radius below, pays the half-log endpoint obstruction.

\begin{theorem}[Gevrey dominance formula]\label{thm:ns:gevrey-dominance}
For the Gevrey-normalized derivative probe
\(\widetilde M_{m,J,t}\), the dominance constant at derivative order
\(m\) has the closed form
\[
  A_m^G(t,J)
  =
  \lambdaJ(t)^{q-\eps}\,
  \frac{m!}{\sigmaJ(t)^m},
  \qquad
  \sigmaJ(t)=\frac{\lambdaJ(t)\rho(t)}{\CG}.
\]
Optimizing over \(m\ge 0\) gives, for a positive constant \(c\),
\[
  \inf_{m\ge 0} A_m^G(t,J)
  \lesssim
  \lambdaJ(t)^{q-\eps}\sqrt{1+\sigmaJ(t)}\,
  e^{-c\sigmaJ(t)} .
\]
Consequently the direct radius-window gate is
\[
  \sigmaJ(t)\ge \log\lambdaJ(t),
\]
up to the harmless normalization constants absorbed into \(c\).  Under
this gate the optimized dominance is smaller than the leading
\(\lambdaJ^{q-\eps}\) scaling by an exponential radius-window factor.
\end{theorem}

\begin{proof}
The \(m\)-th derivative contributes the active-shell factor
\(\lambdaJ(t)^{q-\eps}\) while the Gevrey weight contributes
\(\rho(t)^m/(\CG^m m!)\).  Rewriting through
\(\sigmaJ(t)=\lambdaJ(t)\rho(t)/\CG\) gives the closed form
\(A_m^G(t,J)=\lambdaJ(t)^{q-\eps}m!/\sigmaJ(t)^m\).  Minimising
\(m!/\sigmaJ^m\) by Stirling~\citep{AbramowitzStegun1964}
\((m!\sim m^m e^{-m}\sqrt{2\pi m})\) selects a saddle order \(m\)
comparable to \(\sigmaJ\) (after absorbing Gevrey-normalisation
constants into \(c\), the usual order is \(m\approx\sigmaJ/e\)) and
yields
\[
  \frac{m!}{\sigmaJ^m}
  \lesssim
  \sqrt{1+\sigmaJ}\,e^{-c\sigmaJ},
\]
which is the optimised dominance estimate.  In powers of
\(\lambdaJ\), the exponential decay overtakes the leading
\(\lambdaJ^{q-\eps}\) factor once
\(c\sigmaJ(t)\gtrsim(q-\eps)\log\lambdaJ(t)\); the normalised form is
\(\sigmaJ(t)\ge \log\lambdaJ(t)\).\footnote{Verified in Lean as \texttt{SixBirdsNeedles.NS.GevreyLP.gevreyDominance}; see \Cref{subsec:paper3-per-label-coverage}.}
\end{proof}

\noindent
The half-log analyticity import gives
\(\sigmaJ(t)\gtrsim\log\lambdaJ(t)/\CG\), so the basic gate is met on
the accepted BG-compatible stages.  By
\Cref{thm:ns:no-go-half-log-endpoint}, however, meeting this gate is
not enough for summable predictive propagation: the staged budget has
ratio tending to \(2^{2p}>1\), and
\Cref{thm:ns:no-go-published-analyticity} identifies the corresponding
\(\sqrt{\log\lambdaJ}\) shortfall for the direct contraction route.
The LP-transport rescue below pays this gap by moving the staged
budget into the response space \(H^{-s}\), where the
Littlewood--Paley transport contributes additional decay.

\begin{theorem}[LP-transport rescue]\label{thm:ns:lp-transport-rescue}
Let \(\LPtr:L^2(P_j)\to H^{-s}\) be the LP-transport stage operator with
\(s\ge q-\eps\), and set \(p=q-\eps\).  The half-log staged budget
transports to
\[
  \widetilde B_j(s)
  \simeq
  2^{2(p-s)j}
  \bigl(1+\mu\sqrt{L_j}\bigr)
  e^{-2c\mu\sqrt{L_j}},
  \qquad
  L_j=\log\lambdaJ .
\]
For every \(s>p\),
\[
  \frac{\widetilde B_{j+1}(s)}{\widetilde B_j(s)}
  \longrightarrow
  2^{2(p-s)}<1 .
\]
Consequently \(\sum_j \widetilde B_j(s)<\infty\) on the transported
response space \(H^{-s}\), so the transported residual satisfies the
summable predictive-propagation hypothesis.
\end{theorem}

\begin{proof}
Set \(\widetilde B_j(s)=\|\LPtr\|^2 B_j\), where \(B_j\) is the
half-log budget of NG8 and \(\|\LPtr\|^2\simeq 2^{-2sj}\) is the
audited operator-norm decay from the Littlewood--Paley /
Bony--Chemin--Danchin calculus.  Multiplying through
\(B_j\simeq 2^{2pj}(1+\mu\sqrt{L_j})e^{-2c\mu\sqrt{L_j}}\) gives the
displayed transported budget.  The stage-ratio computation of NG8
applied to \(\widetilde B_j(s)\) leaves the multiplier
\(2^{2(p-s)}\): strictly contractive for \(s>p\), neutral for
\(s=p\), divergent for \(s<p\).\footnote{Verified in Lean as \texttt{SixBirdsNeedles.NS.GevreyLP.lpTransportRescue}; see \Cref{subsec:paper3-per-label-coverage}.}
\end{proof}

\noindent
The lawful LP range \(s\ge q-\eps=p\) is useful in its strictly
transported regime \(s>p\), and \(s\) is taken just above \(p\) to
gain contraction while staying close to the original response scale.
Summability of \(\widetilde B_j(s)\) supplies the summable
predictive-propagation input of \citet{Tsiokos2026NeedleKiller}, which
combined with the predictive membrane on the accepted doubling stages
yields the response-space hypothesis of the layer-dissolving endpoint.

\begin{theorem}[Smooth-data half-log transfer]
\label{thm:ns:smooth-data-half-log-transfer}
On the accepted smooth-data subcarrier \(H^{\gamma_0}\) with
\(\gamma_0>5/2\), the half-log analyticity radius satisfies
\[
  \rho(t)\gtrsim \frac{\log\lambdaJ(t)}{\lambdaJ(t)}
\]
on BG-compatible stages.  This follows by composing the Sobolev/Besov
embedding
\[
  H^{\gamma_0}\hookrightarrow
  \dot B^{-\eps}_{\infty,\infty}
\]
with the half-log analyticity import.  Consequently
\[
  \sigmaJ(t)
  =
  \frac{\lambdaJ(t)\rho(t)}{\CG}
  \gtrsim
  \frac{\log\lambdaJ(t)}{\CG},
\]
so the normalized radius-window gate of
\Cref{thm:ns:gevrey-dominance} is met, and the LP-transport rescue of
\Cref{thm:ns:lp-transport-rescue} gives summable predictive propagation
on \(H^{-s}\) for \(s>q-\eps\).
\end{theorem}

\begin{proof}
Given divergence-free data \(u_0\in H^{\gamma_0}\) with
\(\gamma_0>5/2\), Kato local well-posedness gives a strong solution
\(u(\cdot,t)\in H^{\gamma_0}\) for \(t\in[0,\Tmax(u_0))\).  The
Sobolev/Besov embedding lifts this restart data into
\(H^{\gamma_0}\cap\dot B^{-\eps}_{\infty,\infty}\), and the half-log
analyticity import then supplies the displayed radius bound on
BG-compatible stages.  Rewriting through
\(\sigmaJ(t)=\lambdaJ(t)\rho(t)/\CG\) gives the radius-window input
required by the Gevrey dominance theorem; the gate is met but is not
itself contractive, and \Cref{thm:ns:lp-transport-rescue} supplies the
missing decay on \(H^{-s}\) for \(s>q-\eps\).\footnote{Verified in Lean as \texttt{SixBirdsNeedles.NS.GevreyLP.smoothDataHalfLogTransfer}; see \Cref{subsec:paper3-per-label-coverage}.}
\end{proof}

\noindent
With the amplitude obligation paid, we assemble the records
\(R_1,\ldots,R_9\) in \Cref{sec:main-theorem}.

 

\section{Main Theorem Under Six Birds Closure}
\label{sec:main-theorem}

This section assembles the records \(R_1,\ldots,R_9\) developed in
\Cref{sec:ns-carrier,sec:imports,sec:adequacy-cascade,sec:no-go,sec:gevrey-lp}
into the hypothesis stack consumed by the layer-dissolving membrane
endpoint of \citet{Tsiokos2026NeedleKiller}.  The records and the
typed-context discipline imply \(\Tmax(u_0)=\infty\) on the accepted
smooth-data channel.

\begin{definition}[Six Birds closure records \(R_1\)--\(R_9\)]
\label{def:ns:closure-r1-r9}
The Six Birds closure record bundle for the Navier--Stokes carrier is
the tuple
\[
  \mathsf R=(R_1,R_2,R_3,R_4,R_5,R_6,R_7,R_8,R_9),
\]
where the nine entries are the following proof-carrying records.
\begin{enumerate}[label=$R_\arabic*$.]
\item Adequacy orientation: the carrier data
\((E_{J,t},C_{J,t},\Lphys_{J,t},\DBG_{J,t})\) is ordered so that a
smaller residual \(\Xi\) is better.
\item Exact BG-windowed package: \(\Gamma_j\) preserves
divergence-free structure and the mismatch sector on the packaged
carrier.
\item Null-mode legality: \(\Lphys\) and \(\DBG\) vanish on
\(\Ker C_{J,t}\) as in \Cref{sec:ns-carrier}.
\item Source-admission acceptance: the recurrence with
\(\Omsrc=\Omtr+\Ompr+\Omvisc+\Omtail+\Omstr\) is accepted on the
cascade of \Cref{sec:adequacy-cascade}.
\item BG frequency-doubling staging: the active scales obey
\(\lambda_{j+1}=2\lambda_j\).
\item Dissolving-side LP transport: \(\LPtr:L^2(P_j)\to H^{-s}\) with
\(s\ge q-\eps\) is accepted as the typed transport of
\Cref{sec:gevrey-lp}.
\item Half-log analyticity: restart data in
\(H^{\gamma_0}\cap\dot B^{-\eps}_{\infty,\infty}\), with
\(\gamma_0>3/2-\eps\), carries the half-log lower radius on
BG-compatible stages.
\item Paper-grounded PDE imports: Beale--Kato--Majda continuation,
Bradshaw--Grujic frequency-window continuation, Kato local
well-posedness, half-log analyticity, Littlewood--Paley /
Bony--Chemin--Danchin calculus, and the Sobolev/Besov embedding chain
are accepted at their declared status.
\item Nonclaim boundary: there is no criterion-native replacement of
\(\DBG\), no generic Besov claim, no untyped global comparison, and no
closure outside the smooth-data subcarrier without additional records.
\end{enumerate}
\end{definition}

The bundle records the nine standing hypotheses above as separate
hypothesis components: each \(R_i\) can be stated, checked, and
audited on its own, and \(\mathsf R\) only records that all nine hold
in the same
typed context.  In particular \(\mathsf R\) does not carry
\(\Tmax(u_0)=\infty\); that conclusion is a downstream consequence
obtained by applying the layer-dissolving endpoint through the
typed-context deliverability bridge --- the implication that consumes
the nine \(R_i\)-fields together with the smooth-data hypotheses on
\(u_0\) and concludes \(\Tmax(u_0)=\infty\).

\begin{theorem}[Framework-deliverable Navier--Stokes regularity]
\label{thm:ns:framework-deliverable-regularity}
Under the Six Birds closure record bundle
\(\mathsf R=(R_1,\ldots,R_9)\) of
\Cref{def:ns:closure-r1-r9} and the typed-context discipline of
\citet{Tsiokos2026NeedleKiller} and
\citet{Tsiokos2026XiCompanion}, for every divergence-free initial datum
\(u_0\in H^{\gamma_0}\) with \(\gamma_0>5/2\) on the accepted
smooth-data channel,
\[
  \boxed{\Tmax(u_0)=\infty .}
\]
\end{theorem}

\begin{proof}
The proof is the structural composition described above.
\begin{enumerate}[label=\emph{Step \arabic*.}]
\item The residual calculus is active on the packaged carrier.  The
record \(R_2\) supplies the exact BG-windowed package \(\Gamma_j\) on
\((E_{J,t},C_{J,t},\Lphys_{J,t},\DBG_{J,t})\), and \(R_1\) fixes the
adequacy orientation: smaller \(\Xi\) is better, with order read in
the Loewner cone on the response space.  By \(R_3\), null-mode
legality gives
\(\Lphys_{J,t}(\Ker C_{J,t})=0\) and
\(\DBG_{J,t}(\Ker C_{J,t})=0\), so both families descend to the legal
energy quotient \(E_{C_{J,t}}=E_{J,t}/\Ker C_{J,t}\); on that quotient
the audit form is positive definite by construction, as in
\Cref{thm:ns:aor-quotient-transport-record}.  Therefore the adequacy
residual \(\Xi_{C_{J,t}}(\DBG\mid\Lphys)\) is well-defined and is the
Schur complement
\[
  \begin{aligned}
  \Xi_{C_{J,t}}(\DBG\mid\Lphys)
  &= K_{DD}-K_{DL}K_{LL}^{\dagger}K_{LD},\\
  K_{LL} &= \Lphys C_{J,t}^{\dagger}{\Lphys}^*,\qquad
  K_{DD} = \DBG C_{J,t}^{\dagger}{\DBG}^*,\\
  K_{DL} &= \DBG C_{J,t}^{\dagger}{\Lphys}^* ,
  \end{aligned}
\]
computed on the legal quotient.  The projection identity of
\citet{Tsiokos2026XiCompanion} then rewrites the same residual in
energy-scaled coordinates
\(T_L=\Lphys C_{J,t}^{\dagger/2}\),
\(T_D=\DBG C_{J,t}^{\dagger/2}\) as
\[
  \Xi_{C_{J,t}}(\DBG\mid\Lphys)=T_D(I-P_L)T_D^*,
  \qquad P_L=\operatorname{proj}_{\Ran T_L^*}.
\]
This operator is positive semidefinite, so the orientation of \(R_1\)
is consistent with Loewner order.
\item The source-admission recurrence now propagates the residual along
the Bradshaw--Grujic doubling chain.  The elementary step is
\Cref{thm:ns:source-admission-recurrence}: for one admissible time
step \(h\) on stage \(J\),
\[
  \Xi_{C_{J,t+h}}\bigl(\DBG_{J,t+h}\mid\Lphys_{J,t+h}\bigr)
  \preceqq
  C_{\rm old}\,\mathcal T_h\,
  \Xi_{C_{J,t}}\bigl(\DBG_{J,t}\mid\Lphys_{J,t}\bigr)
  \mathcal T_h^*
  + \Omsrc .
\]
The record \(R_5\) supplies the stage ordering
\(\lambda_{j+1}=2\lambda_j\), so this one-step inequality can be
iterated from stage \(0\) to stage \(J\).  Writing
\(\mathcal T^{(J)}=\mathcal T_{h_{J-1}}\cdots\mathcal T_{h_0}\) and
letting \(\mathcal T^{(J,j)}\) denote the composed transport from
stage \(j+1\) up to stage \(J\), the finite-stage envelope is
\[
  \begin{aligned}
  \Xi_{C_{J,t}}(\DBG_J\mid\Lphys_J)
  &\preceqq
  \Bigl(\prod_{j<J} C_{\rm old}^{(j)}\Bigr)
  \mathcal T^{(J)}
  \Xi_{C_0}(\DBG_0\mid\Lphys_0)
  \bigl(\mathcal T^{(J)}\bigr)^* \\
  &\quad+
  \sum_{j<J}
  \Bigl(\prod_{j<i<J} C_{\rm old}^{(i)}\Bigr)
  \mathcal T^{(J,j)}\,\Omsrc^{(j)}\,
  \bigl(\mathcal T^{(J,j)}\bigr)^* .
  \end{aligned}
\]
Thus all finite-stage growth is accounted for by the transported
initial residual and the accumulated transported source ledgers.  The
adequacy-residual chain rule of \citet{Tsiokos2026XiCompanion} then
refines the accumulated ledger layer by layer: when a prior discharge
witness admits a new native family \(M\) extending \(L\), the chain
rule gives
\[
  \Xi_C(D\mid[L;M])
  \preceqq
  \Xi_C(D\mid L)
  - K_{DM\mid L}K_{MM\mid L}^{\dagger}K_{MD\mid L},
\]
so the extended residual is the previous residual minus a positive
Schur correction.  Applying this contraction successively along
\(\Omsrc\to\Omega_{\rm gap}\to\Omega_{\rm rem}\to
\Omega_{\rm bb}\to\Omega_{\rm mm}\to\Omamp\) peels off each
discharged source component and leaves \(\Omamp\) as the unique
un-discharged ledger at the amplitude layer.
\item The no-go sieve closes every shortcut discharge route exposed by
the cascade of Step~2.  First, \Cref{thm:ns:no-go-scalar-energy}
rules out scalar quadratic budgets such as energy and enstrophy as
Schur probes at any layer: they are not linear native families, so
every cascade discharge must enter through linear native probes.
Next, \Cref{thm:ns:no-go-packing-amplitude,thm:ns:no-go-sparse-shadow}
rule out packing-amplitude and sparseness data as payments at the
\(\Omega_{\rm bb}\) and \(\Omega_{\rm mm}\) bad-cylinder / mismatch
layers unless a sector-matched amplitude bridge is supplied.
\Cref{thm:ns:no-go-marker-only} rules out triad markers and
bad-overlap masks as amplitude-layer discharges without a quantitative
amplitude bridge, while \Cref{thm:ns:no-go-metric-shadow} rules out
geometric-metric reweightings of \(\Omamp\) unless the normalized
metric dominates the fixed BG/BKM amplitude on the carrier scale.
\Cref{thm:ns:no-go-moving-window} excludes moving-support data alone
as a closure record without a finite-termination, fixed-ledger, or
exhaustive-ledger squeeze.  Finally,
\Cref{thm:ns:no-go-half-log-endpoint} rules out the half-log
analyticity lane alone as a discharge of \(\Omamp\), since the staged
budget has ratio tending to \(2^{2p}>1\), and
\Cref{thm:ns:no-go-published-analyticity} records the corresponding
\(\sqrt{\log\lambdaJ}\) shortfall on the direct route; by the same
type discipline as NG2, \Cref{thm:ns:no-go-raw-nonlinearity} rules out
raw nonlinear observables such as fluxes, pressure interactions, and
triad terms as Schur probes.  Thus NG2--NG10 exclude all earlier
terminations of the cascade by public-shadow, marker, metric,
moving-window, endpoint-only, or nonlinear shortcuts.  The
undischarged ledger after the sieve is necessarily \(\Omamp\), whose
payment must come from the linear native probes admitted at the
amplitude layer, namely the analytic records treated in Step~4.
\item The amplitude ledger \(\Omamp\) is paid by the three analytic
records carried by \(R_7\), \(R_6\), and \(R_8\).  First, \(R_8\)
supplies Kato local well-posedness and the Sobolev/Besov embedding
chain: for \(u_0\in H^{\gamma_0}\), \(\gamma_0>5/2\), the strong
solution satisfies \(u(\cdot,t)\in H^{\gamma_0}\) on
\([0,\Tmax(u_0))\), and the embedding lifts the restart data into
\(H^{\gamma_0}\cap\dot B^{-\eps}_{\infty,\infty}\).  The half-log
analyticity record in \(R_7\), composed in
\Cref{thm:ns:smooth-data-half-log-transfer}, then gives on
BG-compatible stages
\[
  \rho(t)\gtrsim \frac{\log\lambdaJ(t)}{\lambdaJ(t)},
  \qquad
  \sigmaJ(t)\geq \log\lambdaJ(t),
\]
after absorbing the harmless Gevrey normalization constant into the
gate constant.  Second, under this gate,
\Cref{thm:ns:gevrey-dominance} gives
\(A_m^G(t,J)=\lambdaJ^{q-\eps}m!/\sigmaJ^m\) and the optimized
dominance
\[
  \inf_{m\geq 0} A_m^G(t,J)
  \lesssim
  \lambdaJ^{q-\eps}\sqrt{1+\sigmaJ}\,e^{-c\sigmaJ}.
\]
Multiplying by the active-shell budget yields the untransported
half-log amplitude budget
\[
  B_j
  =
  C_B\,2^{2pj}\bigl(1+\mu\sqrt{L_j}\bigr)
  e^{-2c\mu\sqrt{L_j}},
  \qquad
  L_j=\log\lambdaJ,\quad p=q-\eps .
\]
By \Cref{thm:ns:no-go-half-log-endpoint}, however,
\(B_{j+1}/B_j\to 2^{2p}>1\), so the bare half-log lane is not
summable on its own.  Third, \(R_6\) supplies the
Littlewood--Paley transport stage operator
\(\LPtr\colon L^2(P_j)\to H^{-s}\), with
\(\|\LPtr\|^2\simeq 2^{-2sj}\) for \(s\geq p\).  The LP-transport
rescue \Cref{thm:ns:lp-transport-rescue} multiplies the half-log
budget by this operator-norm decay and gives the transported budget
on the response space \(H^{-s}\),
\[
  \widetilde B_j(s)
  \simeq
  2^{2(p-s)j}\bigl(1+\mu\sqrt{L_j}\bigr)
  e^{-2c\mu\sqrt{L_j}},
  \qquad
  \frac{\widetilde B_{j+1}(s)}{\widetilde B_j(s)}
  \longrightarrow 2^{2(p-s)}<1
  \quad(s>p).
\]
Thus the transported budget is geometrically contractive on a cofinal
tail of stages and \(\sum_{j\geq0}\widetilde B_j(s)<\infty\).  This
is exactly the summable predictive-propagation budget required by
\citet{Tsiokos2026NeedleKiller}, and it is the input consumed by
Step~5.
\item The summable predictive-propagation budget
\(\sum_j\widetilde B_j(s)<\infty\) on \(H^{-s}\) from Step~4 supplies
the summability input for the layer-dissolving membrane endpoint of
\citet{Tsiokos2026NeedleKiller}.  The remaining endpoint hypotheses are
also furnished by the closure bundle: \(R_2\) and \(R_3\) provide the
exact packaged carrier with null-mode legality activated in Step~1;
\(R_4\), through the source-admission recurrence and the stagewise
budget \(\Theta^Y_J\) of Step~2, supplies the native budget on the
physical-native currency \(K_{LL}\); Step~4 supplies the residual budget
\(\widetilde B_j(s)\preceqq\Xi^Z_J\) for
\(\Xi_C(\DBG\mid\Lphys)\) on the transported response space; and
\(R_1,\ldots,R_8\) supply the all-six transfer discipline.  The
nonclaim record \(R_9\) checks that the endpoint is invoked on the
accepted smooth-data subcarrier, with the readout still \(\DBG\), with
no criterion-native rewrite or generic Besov claim, and with all
comparisons kept in the typed-context disciplines of
\citet{Tsiokos2026NeedleKiller} and
\citet{Tsiokos2026XiCompanion}.  Therefore the endpoint applies to the
transported maps \(S_J=\LPtr\circ\mathcal T^{(J)}\), where
\(\mathcal T^{(J)}\) is the BG-chain transport from Step~2 and \(\LPtr\)
is the LP-transport stage operator from Step~4.  With
\(K^D_J=\DBG_J C_{J,t}^{\dagger}(\DBG_J)^*\), with \(A_J\) the optimal
explanation map at stage \(J\) (the adequacy-promotion map of
\citet{Tsiokos2026XiCompanion}), the endpoint gives the uniform
transported dissolving budget
\[
  S_J K^D_J S_J^*
  \preceqq
  S_J\bigl(A_J\Theta^Y_J A_J^*+\Xi^Z_J\bigr)S_J^*
  \preceqq
  \Theta^D
  \qquad\text{uniformly in }J .
\]
The summability from Step~4 is what makes \(\Theta^D\) uniform rather
than stage-dependent.  This uniform transported dissolving budget is
the activated continuation channel consumed by the Bradshaw--Grujic and
Beale--Kato--Majda continuation criteria in Step~6.
\item The uniform transported dissolving budget from Step~5 now gives
both continuation readings.  For the Bradshaw--Grujic frequency-window
criterion, \citet{BradshawGrujic2017} require uniform control of the
dissolving currency block on the BG-windowed package across the
doubling frequency stages in a left-neighborhood of any proposed finite
terminal time \(T^*\).  The bound
\(S_JK^D_JS_J^*\preceqq\Theta^D\), uniformly in \(J\), is exactly this
hypothesis after transport through the response map \(S_J\); hence the
BG continuation criterion extends the strong solution past \(T^*\).
For the Beale--Kato--Majda criterion, \citet{BealeKatoMajda1984}
continue a strong \(H^{\gamma_0}\) solution past \(T^*\) whenever
\(\int_0^{T^*}\|\omega(\cdot,t)\|_{L^\infty}\,dt<\infty\).  The
Sobolev/Besov chain in \(R_8\), together with the LP-transport rescue
of Step~4, translates the same uniform Loewner bound on \(H^{-s}\) into
uniform pointwise control of the vorticity on each BG window; the BG
window summability and the smooth-data \(H^{\gamma_0}\) embedding then
give the finite vorticity-time integral required by BKM.  The two
criteria therefore read the same activated channel in complementary
forms: Bradshaw--Grujic uses the frequency-window bound, while BKM uses
the classical vorticity-time-integral bound.  If
\(\Tmax(u_0)<\infty\), applying either criterion at
\(T^*=\Tmax(u_0)\) gives continuation strictly past the maximal
lifespan, a contradiction.\footnote{Tracked in the Lean formalization as
\texttt{SixBirdsNeedles.NS.MainTheorem.frameworkDeliverableRegularity}:
a conditional schema over the closure record bundle
\(R_1,\ldots,R_9\) rather than an independent derivation,
which is proved here.  See
\Cref{subsec:paper3-per-label-coverage}.}  Hence
\[
  \Tmax(u_0)=\infty .
\]
\end{enumerate}
\end{proof}

\Cref{sec:ns-aor-instance} recasts \(R_1,\ldots,R_9\) and the
deliverability bridge as the AOR-membership witness for the
smooth-data carrier inside the audit-closure meta-theory of
\citet{Tsiokos2026AOR}, providing an alternative reading of the same
conclusion.

 

\section{NS as an AOR instance}
\label{sec:ns-aor-instance}

\subsection{AOR-membership reading of the NS smooth-data carrier}

\Cref{thm:ns:framework-deliverable-regularity} stated the result in
conditional form: the records \(R_1,\ldots,R_9\), together with the
typed-context bridge, imply \(\Tmax(u_0)=\infty\) for accepted smooth
data.  In this section we recast the same closure stack inside the
audit-closure meta-theory of \citet{Tsiokos2026AOR}.  There an object
belongs to the refinement-stable AOR class when its declared audit
records discharge the structural defects that remain under refinement;
the relevant background is the refinement-stable AOR predicate, the
witness hierarchy \(\mathbf{AOR}^{\mathrm{ref}}\),
\(\mathbf{AOR}^{\mathrm{casc}}\), \(\mathbf{AOR}\) (with each
boundary exhibited by named witness carriers in
\citet{Tsiokos2026AOR}), the eight-stratum structural defect, and
the residual atlas with discharge cascades.  That structural defect
belongs to the AOR
meta-theory and should not be confused with the adequacy residual
\(\Xi_C(D\mid L)\) used earlier.

For the Navier--Stokes carrier the witness data are precisely the
closure records of \Cref{def:ns:closure-r1-r9}: inside the AOR
discipline, those records are the membership witness for the
smooth-data carrier \(\mathsf C_{\mathrm{NS}}\).  Read in this way,
``the carrier does not blow up in reality'' means that the carrier is
audit-closed under the declared refinement-stable AOR presentation;
the closure records are not discarded but reclassified.

\begin{definition}[NS smooth-data carrier as an AOR-instance object]\label{def:ns:aor-instance-object}
The Navier--Stokes smooth-data carrier with initial datum $u_0\in H^{\gamma_0}$, $\gamma_0>5/2$, is the AOR-instance object $\mathsf C_{\mathrm{NS}}$ whose carrier family is the Bradshaw--Grujic frequency-refinement chain $H_{J,t}=E_{J,t}/\Ker C_{J,t}$ indexed by BG stages $J\in\mathbb N$ and time $t<\Tmax(u_0)$, with refinement morphisms $J\mapsto J+1$ of combined kind A3+A6.  Its AOR registries are populated by the records of \S\ref{sec:ns-carrier}--\S\ref{sec:gevrey-lp}: physical-native probes $\Lphys$ and target readout $\DBG$ as observation records; the BG-windowed package $\Gamma_j$, audited transport $\mathcal T_h$, and LP-transport $S_j$ as route records; the source-admission ledger $\Omsrc$ as source records; the BG window $P_{J(t)}$ as the bounded interface; null-mode legality, positivity, and the Loewner recurrence as constraints; the half-log radius bound as a scale record; and the present paper's nonclaim list as the nonclaim register.  The scope $\mathfrak S_{\mathrm{NS}}$ restricts to smooth divergence-free Sobolev data.
\end{definition}

\subsection{Mechanical R-records as membership components}

The first part of the membership witness is mechanical in the precise
sense that the corresponding records are supplied by the carrier
construction, the declared package data, or the approved external
import chain.  These records still matter: they are the typed fields
that make the AOR membership statement well-formed.  They simply do not
carry the substantive source, transport, or half-log discharge work.

\begin{theorem}[Mechanical closure records as AOR-membership components]\label{thm:ns:aor-mechanical-r-records}
Six of the nine standing closure records of $\mathsf C_{\mathrm{NS}}$ admit by-construction or approved-other discharges within scope $\mathfrak S_{\mathrm{NS}}$: $R_1$ (adequacy orientation), $R_2$ (exact BG-windowed package), $R_3$ (null-mode legality), $R_5$ (BG frequency-doubling staging), and $R_9$ (nonclaim boundary) discharge \texttt{by\_construction}; $R_8$ (paper-grounded PDE imports: BKM, BG, Kato, half-log analyticity, LP/BCD, Sobolev/Besov embedding) discharges \texttt{approved\_other} as an external-record citation chain.  These six records contribute the mechanical components of the AOR-membership witness; the substantive components $R_4,R_6,R_7$ follow.
\end{theorem}

\begin{proof}
The six listed records discharge by record-class inspection.
\(R_1\) (adequacy orientation), \(R_2\) (exact BG-windowed
package), and \(R_3\) (null-mode legality) are by-construction
fields of the smooth-data carrier of \Cref{sec:ns-carrier}:
the BG-windowed exact package and the legal-energy quotient
are part of the carrier definition, so the corresponding
discharge witnesses are the constructions themselves.  \(R_5\)
(BG frequency-doubling staging) is by-construction from the
doubling-stage indexing of \Cref{sec:adequacy-cascade}: the
stage map \(J\mapsto J+1\) is the carrier's declared
refinement.  \(R_9\) (nonclaim boundary) is by-construction
from the locked nonclaim list in \Cref{sec:status} together
with the per-paper nonclaim registers.  \(R_8\) (paper-grounded
PDE imports) discharges \texttt{approved\_other} as the
citation chain assembled in \Cref{sec:imports}: each of BKM,
BG, Kato local well-posedness, half-log analyticity,
Littlewood--Paley / Bony--Chemin--Danchin, and Sobolev/Besov
embedding is admitted as a single named external record, and
the bundle is the AOR \texttt{approved\_other} discharge for
\(R_8\).  These six records contribute the mechanical
components of the AOR-membership witness; the substantive
components \(R_4,R_6,R_7\) are treated in the next three
theorems.\footnote{Verified in Lean as \texttt{SixBirdsNeedles.NS.AORInstance.aorMechanicalRRecords}; see \Cref{subsec:paper3-per-label-coverage}.}
\end{proof}

\subsection{Substantive R4, R6, and R7 discharges}

The remaining membership components are substantive because they
connect the source ledger, the quotient transport, and the half-log
analyticity lane to the AOR residual types.  The \(R_4\) record is the
source-admission recurrence of \Cref{sec:adequacy-cascade}; its
transport typing uses the quotient record below.

\begingroup
\hbadness=10000
\begin{theorem}[Substantive discharge of $R_4$ as a membership component]\label{thm:ns:aor-r4-discharge}
The source-admission recurrence record $R_4$ of $\mathsf C_{\mathrm{NS}}$ has primary type $\Delta_{\mathrm{source}}$ and forced secondary types $\Delta_{\mathrm{transport}},\Delta_{\mathrm{scale}},\Delta_{\mathrm{target}},\Delta_{\mathrm{role}}$.  Its discharge status decomposes as $(\Delta_{\mathrm{source}},\texttt{budgeted}\text{ per stage})+(\Delta_{\mathrm{transport}},\texttt{zero})+(\Delta_{\mathrm{scale}},\texttt{cross-cascade to }R_6,R_7)$, with the Loewner recurrence as the budgeted witness and the quotient-transport record \Cref{thm:ns:aor-quotient-transport-record} as the zero-transport witness.  $R_4$ is per-stage budgeted; its asymptotic propagation is not standalone, and is paid by the $R_6/R_7$ cross-cascade.
\end{theorem}
\endgroup

\begin{proof}
The primary type assignment follows from the role of \(R_4\):
it is the source-admission recurrence of
\Cref{thm:ns:source-admission-recurrence}, which by definition
records source obligations on the cascade, so its primary type
is \(\Delta_{\mathrm{source}}\).  The forced-secondary list
\(\Delta_{\mathrm{transport}},\Delta_{\mathrm{scale}},
\Delta_{\mathrm{target}},\Delta_{\mathrm{role}}\) is read off
the AOR forced-residual table for \(\Delta_{\mathrm{source}}\)
\citep{Tsiokos2026AOR}: a source obligation must declare how
the source data are transported, scaled, target-typed, and
role-typed.

The discharge-status decomposition is verified as follows.
\((\Delta_{\mathrm{source}},\texttt{budgeted per stage})\) is
witnessed by the Loewner recurrence of
\Cref{thm:ns:source-admission-recurrence}: the per-stage
inequality
\(\Xi_{C_{J,t+h}}(\DBG_{J,t+h}\mid\Lphys_{J,t+h})
\preceqq C_{\mathrm{old}}\,\mathcal T_h\,
\Xi_{C_{J,t}}(\DBG_{J,t}\mid\Lphys_{J,t})\,\mathcal T_h^*
+\Omsrc\) bounds the next-stage residual by the transported
previous-stage residual plus the source ledger \(\Omsrc\),
which IS the per-stage budget.
\((\Delta_{\mathrm{transport}},\texttt{zero})\) is witnessed
by the quotient-transport record of
\Cref{thm:ns:aor-quotient-transport-record}: the audited
transport \(\mathcal T_h\) descends to the legal energy
quotient, so the transport obligation discharges with no
residual contribution.
\((\Delta_{\mathrm{scale}},\texttt{cross-cascade to
}R_6,R_7)\) is witnessed by the asymptotic statuses supplied
in \Cref{thm:ns:aor-r6-discharge} and
\Cref{thm:ns:aor-r7-discharge}: the per-stage budget alone is
not contractive, so the scale obligation is paid by the
cross-cascade to \(R_6\) (LP-transport contraction) and
\(R_7\) (half-log gate).  This completes the decomposition.\footnote{Verified in Lean as \texttt{SixBirdsNeedles.NS.AORInstance.aorR4Discharge}; see \Cref{subsec:paper3-per-label-coverage}.}
\end{proof}

The quotient-transport record is the carrier-space bridge consumed by
the \(R_4\) discharge.  It packages the fact that audited transport is
well typed on the legal quotient, the same quotient on which the
Schur-complement recurrence is read.

\begin{theorem}[Quotient-transport record $R_{4q}$]\label{thm:ns:aor-quotient-transport-record}
The audited transport map $\mathcal T_h$ of $\mathsf C_{\mathrm{NS}}$ descends to the legal energy quotient:
\[
  \mathcal T_h(\Ker C_{J,t}) \;\subseteq\; \Ker C_{J,t+h}
\]
for every admissible $(J,t,h)$, so the induced map $\bar{\mathcal T}_h:E_{J,t}/\Ker C_{J,t}\to E_{J,t+h}/\Ker C_{J,t+h}$ is well-defined and the Loewner recurrence is typed on the legal quotient.
\end{theorem}

\begin{proof}
On the smooth-data Sobolev carrier \(H^{\gamma_0}\) with
\(\gamma_0>5/2\), the audit form \(C_{J,t}\) is
positive-definite on \(E_{J,t}\), so \(\Ker C_{J,t}=\{0\}\)
and the descent is automatic: the inclusion
\(\mathcal T_h(\Ker C_{J,t})\subseteq\Ker C_{J,t+h}\) holds
vacuously, and the induced map
\(\bar{\mathcal T}_h:E_{J,t}/\Ker C_{J,t}\to
E_{J,t+h}/\Ker C_{J,t+h}\) is the transport map itself.  Hence
the Loewner recurrence of
\Cref{thm:ns:source-admission-recurrence} is well-typed on
the legal quotient.\footnote{Verified in Lean as \texttt{SixBirdsNeedles.NS.AORInstance.aorQuotientTransportRecord}; see \Cref{subsec:paper3-per-label-coverage}.}
\end{proof}

The next two substantive discharges are the analytic tail of the cascade:
\(R_6\) supplies LP-transport contraction, and \(R_7\) supplies the
smooth-data half-log transfer.

\begin{theorem}[Substantive discharge of $R_6$ as a membership component]\label{thm:ns:aor-r6-discharge}
The dissolving-side LP-transport record $R_6$ of $\mathsf C_{\mathrm{NS}}$ has primary type $\Delta_{\mathrm{transport}}$ and forced secondary types $\Delta_{\mathrm{scale}},\Delta_{\mathrm{limit}},\Delta_{\mathrm{target}},\Delta_{\mathrm{role}},\Delta_{\mathrm{presentation}}$.  Its discharge status decomposes as $(\Delta_{\mathrm{transport}},\texttt{zero})+(\Delta_{\mathrm{scale}},\texttt{contractive on a cofinal tail})+(\Delta_{\mathrm{limit}},\texttt{summable})$ for every $s>p=q-\eps$, with the LP-transport rescue \Cref{thm:ns:lp-transport-rescue} as the witness for both the contractive-tail and summable statuses; the stage ratio satisfies $\widetilde B_{j+1}(s)/\widetilde B_j(s)\to 2^{2(p-s)}<1$.
\end{theorem}

\begin{proof}
The primary type assignment follows from the role of
\(R_6\): it is the LP-transport rescue of
\Cref{thm:ns:lp-transport-rescue}, which carries the
dissolving-side response through the audited transport map
\(\LPtr\colon L^2(P_j)\to H^{-s}\), so its primary type is
\(\Delta_{\mathrm{transport}}\).  The forced-secondary list
follows from the AOR forced-residual table for
\(\Delta_{\mathrm{transport}}\) \citep{Tsiokos2026AOR}: a
transport obligation must declare scale, limit, target, role,
and presentation typings.

The discharge-status decomposition follows from
\Cref{thm:ns:lp-transport-rescue}.
\((\Delta_{\mathrm{transport}},\texttt{zero})\) is witnessed
because the transport map \(\LPtr\) IS the audited transport;
the transport obligation is closed by exhibiting the map.
\((\Delta_{\mathrm{scale}},\texttt{contractive on a cofinal
tail})\) is witnessed by the stage-ratio limit
\(\widetilde B_{j+1}(s)/\widetilde B_j(s)\to 2^{2(p-s)}<1\)
for \(s>p=q-\eps\) proved in
\Cref{thm:ns:lp-transport-rescue}: the transported budget
strictly contracts on a cofinal tail of stages.
\((\Delta_{\mathrm{limit}},\texttt{summable})\) is witnessed
by \(\sum_j\widetilde B_j(s)<\infty\), the summability
consequence of the same stage-ratio bound.\footnote{Verified in Lean as \texttt{SixBirdsNeedles.NS.AORInstance.aorR6Discharge}; see \Cref{subsec:paper3-per-label-coverage}.}
\end{proof}

\begin{theorem}[Substantive discharge of $R_7$ as a membership component]\label{thm:ns:aor-r7-discharge}
The half-log analyticity / smooth-data transfer record $R_7$ of $\mathsf C_{\mathrm{NS}}$ has primary type $\Delta_{\mathrm{scale}}$ and forced secondary types $\Delta_{\mathrm{limit}},\Delta_{\mathrm{horizon}},\Delta_{\mathrm{source}},\Delta_{\mathrm{target}},\Delta_{\mathrm{presentation}}$.  Its discharge status decomposes as $(\Delta_{\mathrm{scale}},\texttt{zero})+(\Delta_{\mathrm{limit}},\texttt{asymptotic\_budgeted, basic gate})+(\Delta_{\mathrm{horizon}},\texttt{out\_of\_scope/nonclaim})$, with \Cref{thm:ns:smooth-data-half-log-transfer} as the zero-scale witness, the half-log radius bound $\rho(t)\gtrsim\log\lambdaJ(t)/\lambdaJ(t)$ as the basic-gate witness, and the smooth-data Sobolev restriction as the horizon nonclaim.  $R_7$ is not contractive on its own; contraction is supplied by $R_6$.
\end{theorem}

\begin{proof}
The primary type assignment follows from the role of
\(R_7\): it is the smooth-data half-log transfer of
\Cref{thm:ns:smooth-data-half-log-transfer}, which supplies
the half-log radius bound and hence the radius-window gate,
so its primary type is \(\Delta_{\mathrm{scale}}\).  The
forced-secondary list follows from the AOR forced-residual
table for \(\Delta_{\mathrm{scale}}\) \citep{Tsiokos2026AOR}:
a scale obligation must declare limit, horizon, source,
target, and presentation typings.

The discharge-status decomposition is verified as follows.
\((\Delta_{\mathrm{scale}},\texttt{zero})\) is witnessed by
\Cref{thm:ns:smooth-data-half-log-transfer}: the half-log
radius bound \(\rho(t)\gtrsim \log\lambdaJ(t)/\lambdaJ(t)\)
supplies the scale obligation directly.
\((\Delta_{\mathrm{limit}},\texttt{asymptotic\_budgeted,
basic gate})\) is witnessed by the gate
\(\sigmaJ(t)\geq \log\lambdaJ(t)\) derived in
\Cref{thm:ns:smooth-data-half-log-transfer}: the gate
supplies a uniform asymptotic budget but is not strictly
contractive on its own.
\((\Delta_{\mathrm{horizon}},\texttt{out\_of\_scope/nonclaim})\)
is witnessed by the smooth-data Sobolev restriction
\(u_0\in H^{\gamma_0}\), \(\gamma_0>5/2\): every horizon
beyond this subcarrier is recorded as a nonclaim in
\Cref{nonclaim:ns:aor-instance-smooth-data-only}.  The
contractive content for \(R_7\) is supplied by \(R_6\)
through the cross-cascade.\footnote{Verified in Lean as \texttt{SixBirdsNeedles.NS.AORInstance.aorR7Discharge}; see \Cref{subsec:paper3-per-label-coverage}.}
\end{proof}

\subsection{Statement of the AOR-instance theorem}

The preceding records now form a joint discharge package.  The
conditional theorem \Cref{thm:ns:framework-deliverable-regularity}
reads the same package as hypotheses implying regularity.  The
AOR-instance theorem reads the package as the membership witness for
\(\mathsf C_{\mathrm{NS}}\) in the refinement-stable AOR class.

\begin{theorem}[NS smooth-data carrier in $\mathbf{AOR}^{\mathrm{ref}}$]\label{thm:ns:in-refstableaor}
Under the discharges of \Cref{thm:ns:aor-mechanical-r-records}, \Cref{thm:ns:aor-r4-discharge}, \Cref{thm:ns:aor-r6-discharge}, and \Cref{thm:ns:aor-r7-discharge}, together with the quotient-transport record of \Cref{thm:ns:aor-quotient-transport-record}, the Navier--Stokes smooth-data carrier $\mathsf C_{\mathrm{NS}}$ belongs to $\mathbf{AOR}^{\mathrm{ref}}_{\mathfrak S_{\mathrm{NS}}}$.  By the AOR meta-theory's discipline, audit-closure membership is the record-level reading of ``the carrier does not blow up in reality''; the framework-deliverable regularity statement \Cref{thm:ns:framework-deliverable-regularity} is the same content read in conditional form, with the closure stack presented as hypotheses rather than as the membership witness.
\end{theorem}

\begin{proof}
We verify the refinement-stable AOR membership predicate of
\citet{Tsiokos2026AOR} for \(\mathsf C_{\mathrm{NS}}\) using
the five discharge records assembled above.

Refinement-stable AOR membership requires (i) cascade
stability --- every \(\Sigma_{\mathrm{residual}}\)-atom of \(\mathsf C_{\mathrm{NS}}\)
admits a discharge witness whose cost-table cascade reaches the
empty record in finitely many steps using the base AOR-8
statuses --- and (ii) closure under refinement chains --- every
residual atom that survives only at the directed colimit of a
refinement chain into \(\mathsf C_{\mathrm{NS}}\) is paid by
one of the four asymptotic statuses
\(\texttt{eventual\_zero}\), \(\texttt{contractive}\),
\(\texttt{summable}\), \(\texttt{asymptotic\_budgeted}\)
\citep{Tsiokos2026AOR}.

Cascade stability: \Cref{thm:ns:aor-mechanical-r-records}
shows that \(R_1,R_2,R_3,R_5,R_8,R_9\) discharge by
construction or as approved external imports, so the
corresponding residual atoms contribute terminal statuses to
the cost-table cascade without further obligation.
\Cref{thm:ns:aor-r4-discharge} discharges the \(R_4\)
source-admission residual as \texttt{budgeted per stage}
through the Loewner recurrence of
\Cref{thm:ns:source-admission-recurrence}, together with the
quotient-transport record of
\Cref{thm:ns:aor-quotient-transport-record} closing the
\(\Delta_{\mathrm{transport}}\) component.  The \(R_4\)
\(\Delta_{\mathrm{scale}}\) component is paid by the
cross-cascade to \(R_6\) and \(R_7\), which by
\Cref{thm:ns:aor-r6-discharge} and
\Cref{thm:ns:aor-r7-discharge} supply the
\texttt{contractive}/\texttt{summable} and basic-gate
statuses respectively.  Composing these per-record discharges
through the AOR cost-table walk of \citet{Tsiokos2026AOR}
gives, for every \(\Sigma_{\mathrm{residual}}\)-atom of \(\mathsf C_{\mathrm{NS}}\),
a finite discharge cascade reaching the empty record.

Refinement closure: the BG frequency-refinement chain
\(\{H_{J,t}\}_{J\in\mathbb N}\) is the carrier's declared
refinement diagram (kind \(A_3+A_6\) per
\Cref{def:ns:aor-instance-object}).  Limit-stage residuals
along this chain occur only for the
\(\Delta_{\mathrm{limit}}\) component of \(R_6\); by
\Cref{thm:ns:aor-r6-discharge} that component is statused
\texttt{summable}.  Limit-stage residuals for the
\(\Delta_{\mathrm{horizon}}\) component of \(R_7\) are
explicitly classified \texttt{out\_of\_scope/nonclaim} by
\Cref{thm:ns:aor-r7-discharge} and by the nonclaim of
\Cref{nonclaim:ns:aor-instance-smooth-data-only}.  No other
residual stratum survives to the colimit.

Therefore \(\mathsf C_{\mathrm{NS}}\) satisfies both
predicates of refinement-stable AOR membership, and
\(\mathsf C_{\mathrm{NS}}\in
\mathbf{AOR}^{\mathrm{ref}}_{\mathfrak S_{\mathrm{NS}}}\).\footnote{Tracked in the Lean formalization as
\texttt{SixBirdsNeedles.NS.AORInstance.inRefStableAOR}: a
conditional schema over the carrier-attached discharge
assignment and the per-residual-primary-type dispatch rather
than an independent derivation, which is proved here.  See
\Cref{subsec:paper3-per-label-coverage}.}
\end{proof}
 

\section{Status, Limitations, and External Interpretation}
\label{sec:status}

The principal result is
\Cref{thm:ns:framework-deliverable-regularity}: given the closure
records \(R_1,\ldots,R_9\), every divergence-free initial datum
\(u_0\in H^{\gamma_0}\) with \(\gamma_0>5/2\) on the accepted
smooth-data channel satisfies \(\Tmax(u_0)=\infty\).  The companion
\Cref{thm:ns:in-refstableaor} of \Cref{sec:ns-aor-instance} recasts
the same records as the AOR-membership witness for the smooth-data
carrier \(\mathsf C_{\mathrm{NS}}\), so that under the
refinement-stable audit-closure discipline of
\citet{Tsiokos2026AOR} membership reads as the record-level statement
that the carrier does not blow up in reality.  Read together or
separately, the manuscript is not a Clay--Millennium claim.  The
owned contribution is the Navier--Stokes specialisation of these
frameworks:
\begin{itemize}
\item the smooth-data carrier, the fixed dissolving readout, the
physical-native probe family, the BG-windowed exact package, and
null-mode legality (\Cref{sec:ns-carrier});
\item the source-admission ledger and the stage-to-stage recurrence
(\Cref{sec:adequacy-cascade});
\item the no-go family NG1--NG10 excluding shortcut payments
(\Cref{sec:imports,sec:no-go});
\item the Gevrey radius-window branch, LP-transport rescue, and
smooth-data half-log transfer (\Cref{sec:gevrey-lp});
\item the closure record bundle and the conditional-form statement
of \Cref{thm:ns:framework-deliverable-regularity}
(\Cref{sec:main-theorem});
\item the AOR-instance recasting of the same records and its
partial-status disclosure (\Cref{sec:ns-aor-instance}).
\end{itemize}

The residual calculus is imported from
\citet{Tsiokos2026XiCompanion}: the Schur-complement object
\(\Xi_C(D\mid L)\), the projection identity, the optimal residual
identity, exact adequacy, adequacy promotion, the chain rule, the
data-processing identities for post-processing and native coarsening,
and the blind-spot witness characterisation, together with the
typed-context audit discipline.

The endpoint machinery is imported from
\citet{Tsiokos2026NeedleKiller}: the layer-dissolving membrane
endpoint theorem, the all-six discipline, predictive membranes, the
summable predictive-propagation lemma, and the predictive-native
membrane construction.

The audit-closure meta-theory is imported from
\citet{Tsiokos2026AOR}: the refinement-stable AOR class and the rule
that AOR membership is read as audit-closed carrier behavior under the
declared record discipline.  This paper does not redevelop that
meta-theory; it supplies the Navier--Stokes membership witness by
presenting \(R_1,\ldots,R_9\) as the carrier's AOR records.

The PDE inputs --- Beale--Kato--Majda continuation,
Bradshaw--Grujic frequency-window continuation, Kato local
well-posedness, Foias--Temam / Herbst--Skibsted / Wang half-log
analyticity, Littlewood--Paley / Bony--Chemin--Danchin calculus,
and the Sobolev/Besov embedding chain --- enter through the import
record \(R_8\), with conclusions supplied by the cited PDE literature
and consumed under named hypotheses.  We compose them through the
audited residual cascade and no-go sieve into the
conditional-form T10; the PDE results themselves are neither
extended nor rederived here.

\begin{nonclaim}[Outside-framework reading]\label{nonclaim:ns:outside-framework}
Outside the Six Birds closure framework, \Cref{thm:ns:framework-deliverable-regularity} reads as the conditional implication
\[
  \bigl(\mathsf R\text{ holds on the smooth-data carrier}\bigr)
  \Longrightarrow
  \bigl(\Tmax(u_0)=\infty\bigr).
\]
This paper does not claim:
\begin{enumerate}
\item a standalone Clay-Millennium
solution~\citep{Fefferman2000NavierStokes} independent of the Six
Birds closure discipline;
\item generic critical-Besov data regularity outside the smooth-data
subcarrier \(H^{\gamma_0}\), \(\gamma_0>5/2\), in the sense of
critical-space well-posedness
results~\citep{KochTataru2001,Cannone1995,LemarieRieusset2023};
\item a criterion-native rewriting of \(\DBG\) in place of
physical-native adequacy, as ruled out by
\Cref{thm:ns:no-go-criterion-native};
\item that any single PDE input alone closes the problem, since the
argument uses the full \(R_1,\ldots,R_9\) bundle architecture;
\item engagement with Caffarelli--Kohn--Nirenberg partial
regularity~\citep{CaffarelliKohnNirenberg1982}, which concerns
Leray--Hopf weak solutions~\citep{Leray1934,Hopf1951} on a different
carrier class;
\item an untyped global comparison, since every comparison in the proof
is made inside the typed-context discipline;
\item a full-log analyticity radius theorem, since the paper uses only
the available half-log analyticity import together with the
LP-transport rescue.
\end{enumerate}
\end{nonclaim}

\begin{nonclaim}[AOR-instance reading restricted to smooth-data subcarrier]\label{nonclaim:ns:aor-instance-smooth-data-only}
The AOR-instance reading of \Cref{sec:ns-aor-instance}, including \Cref{thm:ns:in-refstableaor}, is restricted to the smooth-data Sobolev subcarrier $u_0\in H^{\gamma_0}$, $\gamma_0>5/2$.  No AOR-membership claim is made for generic critical-Besov data, weak solutions, distributional carriers, or any class outside the declared scope $\mathfrak S_{\mathrm{NS}}$.  The $\Delta_{\mathrm{horizon}}$ residual stratum is explicitly statused as \texttt{out\_of\_scope/nonclaim} in \Cref{thm:ns:aor-r7-discharge}.
\end{nonclaim}

\begin{nonclaim}[No independent Clay-Millennium claim via AOR instance]\label{nonclaim:ns:aor-instance-no-clay-claim}
\Cref{thm:ns:in-refstableaor} establishes membership of $\mathsf C_{\mathrm{NS}}$ in $\mathbf{AOR}^{\mathrm{ref}}$ under the discharges of $R_1$--$R_9$ and the carrier-attached discharge assignment.  Independently of those audit records and of the six declared PDE imports of $R_8$, the result is not a Clay-Millennium claim; the AOR-instance reading reframes the closure stack as the membership witness but does not assert global Navier--Stokes regularity outside the audit-record presentation.
\end{nonclaim}

The nonclaim boundary is the same under both readings: the result
asserts only that if the antecedent \(\mathsf R=(R_1,\ldots,R_9)\) is
formed on the accepted smooth-data carrier, then the maximal
strong-solution lifespan is infinite.  The AOR-instance recasting
presents the same antecedent as a membership-witness package; it does
not claim PDE regularity independent of audit records, does not extend
to carriers beyond the smooth-data subcarrier, and does not yet claim
data-bearing R-records.

Future work may sharpen the bundle --- for example, by replacing the
combined half-log analyticity and LP-transport rescue with a stronger
analyticity radius that delivers contraction on its own, or by
promoting the R-record acceptance fields from \(\mathrm{Prop}\)-valued
predicates to data-bearing structures.  Those strengthenings would
change the external obligations, but they lie beyond the present scope.

 
\appendix


\section{NS-1--NS-61 Step Coverage}
\label{app:step-coverage}

This appendix maps the underlying Navier--Stokes construction track,
steps NS-1 through NS-61, onto the paper body.  The table groups the
steps by construction phase rather than listing them individually;
each row names the phase, its step range, and the body section in
which the corresponding mathematical content is stated, proved, or
fenced as a nonclaim.

\begin{table}[ht]
\centering
\footnotesize
\begin{tabularx}{\textwidth}{@{}>{\raggedright\arraybackslash}X
                                l
                                >{\raggedright\arraybackslash}X@{}}
\toprule
Phase & Step range & Body section \\
\midrule
Phase 1: NS carrier and BG/BKM readout
  & NS-1--NS-5
  & Carrier setup, \Cref{sec:ns-carrier} \\
Phase 2: Literature imports and criterion-native gap
  & NS-6--NS-10
  & Imports and NG1, \Cref{sec:imports} \\
Phase 3: Source-admission ledger and recurrence
  & NS-11--NS-16
  & Adequacy cascade, \Cref{sec:adequacy-cascade} \\
Phase 4: Residual cascade refinements
  & NS-17--NS-27
  & Cascade bridge, \Cref{sec:adequacy-cascade} \\
Phase 5: No-go theorems NG2--NG10
  & NS-28--NS-43
  & No-go sieve, \Cref{sec:no-go} \\
Phase 6: Gevrey radius-window probes and gate
  & NS-44--NS-49
  & Gevrey probes and dominance, \Cref{sec:gevrey-lp} \\
Phase 7: Half-log endpoint obstruction
  & NS-50--NS-53
  & NG8 and callbacks, \Cref{thm:ns:no-go-half-log-endpoint,sec:gevrey-lp} \\
Phase 8: LP-transport rescue and smooth-data transfer
  & NS-54--NS-60
  & LP rescue and transfer, \Cref{sec:gevrey-lp} \\
Phase 9: Main theorem and closure records
  & NS-61
  & Closure bundle and T10, \Cref{sec:main-theorem} \\
\bottomrule
\end{tabularx}
\caption{Construction-step coverage by phase.}
\label{tab:ns-step-coverage}
\end{table}

Per-step entries live in the project's traceability and inventory
artifacts rather than in this typeset table.  The operational records
sit under \path{formalization/traceability/} alongside the source-items
inventory and Lean manifest cross-references; Appendix~D records the
formal artifact map.

 \section{Bridge Atlas and Import Statuses}
\label{app:bridge-atlas}

This appendix records the two audit maps used by the body.  The first
table lists the cross-context bridges by which the residual calculus,
the layer-dissolving endpoint, the Navier--Stokes carrier, and the PDE
imports are connected.  The second table records the declared status of
the six PDE imports: each entry is consumed under named hypotheses, and
the analytic conclusion is supplied by the corresponding classical
result rather than reproved here.

\begin{table}[ht]
\centering
\footnotesize
\begin{tabularx}{\textwidth}{@{}>{\raggedright\arraybackslash}p{0.24\textwidth}
                                >{\raggedright\arraybackslash}p{0.24\textwidth}
                                >{\raggedright\arraybackslash}X@{}}
\toprule
Source context & Target context & Bridge content \\
\midrule
\citet{Tsiokos2026XiCompanion}
  & NS carrier (this paper)
  & The adequacy residual \(\Xi_C(D \mid L)\) is specialised to
    \(\Xi_{C_{J,t}}(\DBG \mid \Lphys)\); native coarsening,
    promotion, and blind-spot witness statements apply at the NS
    instantiation. \\
\citet{Tsiokos2026NeedleKiller}
  & NS carrier
  & The layer-dissolving membrane endpoint, summable predictive
    propagation, the predictive-native membrane construction, and the
    all-six discipline apply to \(R_1,\ldots,R_9\). \\
NS carrier
  & BG continuation channel
  & Uniform boundedness of \(\DBG_{J,t}\) on doubling stages converts
    to global continuation through the Bradshaw--Grujic
    frequency-window criterion. \\
NS carrier
  & BKM continuation channel
  & The vorticity-time integral feeds the Beale--Kato--Majda
    continuation criterion. \\
Smooth-data subcarrier \(H^{\gamma_0}\)
  & Half-log restart hypothesis space
  & The Sobolev/Besov embedding chain lifts smooth-data restart data
    into \(H^{\gamma_0} \cap \dot B^{-\eps}_{\infty,\infty}\). \\
Source-admission ledger \(\Omsrc\)
  & Schur-block bookkeeping
  & The five-component decomposition is tied to the Navier--Stokes
    equation terms through the Littlewood--Paley /
    Bony--Chemin--Danchin calculus and the Helmholtz / Leray
    projection. \\
\bottomrule
\end{tabularx}
\caption{Cross-paper and cross-context bridges used by the NS argument.}
\label{tab:ns-bridge-atlas}
\end{table}

The bridge atlas shows how the formal comparison is transported across
contexts.  The import-status table records the PDE side of the same
audit: the cited form names the classical result, the conclusion column
states what the proof-carrying record supplies under named hypotheses,
and the final column points to the body sections where the record is
invoked.

\begin{table}[ht]
\centering
\footnotesize
\begin{tabularx}{\textwidth}{@{}>{\raggedright\arraybackslash}p{0.19\textwidth}
                                >{\raggedright\arraybackslash}p{0.23\textwidth}
                                >{\raggedright\arraybackslash}X
                                >{\raggedright\arraybackslash}p{0.16\textwidth}@{}}
\toprule
Import & Cited form & Named-hypothesis conclusion & Body section invoked in \\
\midrule
Beale--Kato--Majda
  & BKM continuation criterion
  & \(\int_0^T \|\omega(\cdot,t)\|_{L^\infty}\,dt < \infty
     \Rightarrow \Tmax > T\).
  & \Cref{sec:imports} \\
Bradshaw--Grujic
  & BG frequency-window criterion
  & Uniform boundedness of \(\DBG\) on doubling stages with
    \(\lambda_{j+1}=2\lambda_j\) implies \(\Tmax=\infty\).
  & \Cref{sec:imports} \\
Kato
  & Local well-posedness for divergence-free \(H^{\gamma_0}\) data,
    \(\gamma_0>5/2\)
  & Unique strong solution with maximal lifespan
    \(\Tmax \in (0,\infty]\).
  & \Cref{sec:imports} \\
Half-log analyticity
  & Foias--Temam / Herbst--Skibsted / Wang
  & \(\rho(t) \gtrsim \log\lambda_J / \lambda_J\) on BG-compatible
    stages for restart data in
    \(H^{\gamma_0} \cap \dot B^{-\eps}_{\infty,\infty}\),
    \(\gamma_0>3/2-\eps\).
  & \Cref{sec:imports,sec:gevrey-lp} \\
Littlewood--Paley / Bony--Chemin--Danchin
  & Paradifferential calculus
  & Shell orthogonality, commutator estimates, and operator-norm
    bounds for smooth multipliers.
  & \Cref{sec:imports,sec:adequacy-cascade,sec:gevrey-lp} \\
Sobolev/Besov embedding
  & Embedding chain
  & \(H^{\gamma_0} \hookrightarrow
    \dot B^{-\eps}_{\infty,\infty}\) for \(\gamma_0>3/2\) and
    admissible \(\eps>0\).
  & \Cref{sec:imports} \\
\bottomrule
\end{tabularx}
\caption{Declared status of the PDE imports consumed by the proof.}
\label{tab:ns-import-status}
\end{table}

 \section{Residual Tree and No-Go Proofs}
\label{app:residual-tree}

This appendix records the residual cascade used in
\Cref{sec:adequacy-cascade} and gives appendix-grade proofs of the
schema no-gos stated in \Cref{sec:no-go}.  Each row of the cascade
table names the obstruction node, says what it accounts for, and
points to the import, refinement, no-go, or rescue that addresses
it.

\begin{table}[ht]
\centering
\footnotesize
\begin{tabularx}{\textwidth}{@{}>{\raggedright\arraybackslash}p{0.26\textwidth}
                                >{\raggedright\arraybackslash}X
                                >{\raggedright\arraybackslash}p{0.28\textwidth}@{}}
\toprule
Node & Accounts for & Resolved by \\
\midrule
\(\Xi_{C_{J,t}}(\DBG \mid \Lphys)\)
  & Physical-native adequacy residual.
  & Source-admission recurrence T4 and the Gevrey/LP rescue. \\
\(\Omsrc\)
  & New contributions per step: transport, pressure, viscosity,
    tail, and stretching.
  & The ledger decomposition of \Cref{sec:adequacy-cascade}. \\
\(\Omega_{\rm gap}\)
  & Gap after the BKM, Bradshaw--Grujic, and Kato imports.
  & The imports of \Cref{sec:imports} and the recurrence
    \Cref{thm:ns:source-admission-recurrence}. \\
\(\Omega_{\rm rem}\)
  & Remainder after half-log analyticity, Littlewood--Paley /
    Bony--Chemin--Danchin calculus, and Sobolev/Besov embedding.
  & The analytic imports of \Cref{sec:imports}. \\
\(\Omega_{\rm bb}\)
  & Bad-set / bad-block obstruction.
  & NG3--NG6 in \Cref{sec:no-go}. \\
\(\Omega_{\rm mm}\)
  & Sector-mismatch obstruction.
  & NG5--NG6 in \Cref{sec:no-go}. \\
\(\Omamp\)
  & Amplitude obstruction remaining after the no-go sieve.
  & The Gevrey radius-window branch and LP transport in
    \Cref{sec:gevrey-lp}. \\
\(\Omend\)
  & Endpoint obstruction at the end of the cascade.
  & NG8 / NG9 callbacks and the LP-transport rescue
    \Cref{thm:ns:lp-transport-rescue}. \\
\bottomrule
\end{tabularx}
\caption{Residual cascade nodes and their assigned resolutions.}
\label{tab:ns-residual-tree}
\end{table}

The proofs below cover the no-gos whose body presentations were only
sketched; NG1 is witnessed in \Cref{sec:imports}, NG2 is a body-level
type argument, and NG8 contains its stage-ratio computation in the
body itself.

\begin{proof}[Proof of \Cref{thm:ns:no-go-packing-amplitude}]
Let a bad set at scale \(r\) have measure comparable to \(r^{3-d}\)
and let the relevant windowed amplitude be \(r^{-\alpha}\).  The
contribution to the squared BG readout is then the product of size and
squared amplitude,
\[
  \|\DBG u\|^2_{\text{bad scale}}
  \sim r^{3-d} r^{-2\alpha}
  = r^{(3-d)-2\alpha}.
\]
Taking the square root gives the BG-amplitude contribution
\(r^{(3-d)/2-\alpha}\).  Packing improves the estimate only when this
power is positive, equivalently \(\alpha<(3-d)/2\).  If
\(\alpha\ge (3-d)/2\), the small-set gain is exactly canceled or worse
by amplitude growth, so packing alone cannot dominate \(\Omamp\).  A
lawful payment must therefore add sector matching, tail smallness, or
another amplitude bridge recorded in the cascade.
\end{proof}

\begin{proof}[Proof of \Cref{thm:ns:no-go-sparse-shadow}]
Sparse data controls mass inside the sparse field's declared support,
but the Schur block \(\Lphys C_{J,t}^{\dagger}(\Lphys)^*\) also sees
interactions with the surrounding non-sparse region.  The bad-cylinder /
bad-geometry ledger \(\Omega_{\rm bb}\) is precisely the place where
those interactions are audited.  If the sparse field is not contained
in the same bad-bad or mismatch sector, its estimate pays only the
sparse subset and leaves the cross-sector part unpriced.  The missing
part is not a defect in sparseness; it is a type mismatch between a
set-local estimate and a sector-resolved Schur block.  Hence the proof
splits: either state sector containment, or supply a positive mismatch
record \(\Omega_{\rm mm}\succeq 0\).  Without one of these additions,
sparse-shadow data cannot pay \(\Omega_{\rm bb}\).
\end{proof}

\begin{proof}[Proof of \Cref{thm:ns:no-go-marker-only}]
A marker field has a discrete codomain, typically \(\{0,1\}\), and
records only whether a defect is present or which class it belongs to.
If such a marker is inserted as a probe, the associated Gram block
\(K_{LL}\) is determined by the marker values and the audit weights;
it does not scale with the amplitude of the underlying velocity,
vorticity, or triad contribution.  Thus two carrier states with the
same marker pattern but different BG amplitude produce the same marker
Schur data while requiring different payments in \(\Omamp\).  The
marker can locate the obstruction, but it cannot bound its size.  A
valid promotion therefore needs a quantitative bridge: either a new
linear native probe admitted through the source-admission recurrence,
or a record-grade \(M^{\rm bb}\) bound that carries the missing
amplitude information.
\end{proof}

\begin{proof}[Proof of \Cref{thm:ns:no-go-metric-shadow}]
Changing the carrier metric can shrink the apparent residual by
changing the scale on which the residual is measured.  If a
Gevrey-flavored weight or a Caffarelli--Kohn--Nirenberg-style parabolic
metric~\citep[and the later refinements
of][]{Lin1998CKN,Vasseur2007PartialReg} grows on the bad region, then
the weighted residual there may become small simply because the metric
has divided by a large weight.  The fixed BG observable \(\DBG\),
however, is not changed by this renormalization, and the continuation
channel reads its normalized amplitude.  Constantin--Fefferman
vorticity-direction geometry has the same status for this purpose: it
may organize geometry, but it is not automatically a dominance
estimate for \(\DBG\).  To become a payment, the metric input must
state an explicit dominance constant on the normalized BG scale and a
time gate preventing the weight from degenerating on the active stage.
\end{proof}

\begin{proof}[Proof of \Cref{thm:ns:no-go-moving-window}]
A moving-window or derivative-stack record identifies where the active
frequency content is supported, for instance through windows containing
\(\nabla^mP_Ju(\cdot,t)\) for several \(m\).  Support information is
not a budget: it does not say that the support stops moving, stays
inside a fixed ledger, or is exhausted by already audited ledgers.  If
there is no finite-termination record, the support may keep moving
through infinitely many stages.  If there is no fixed-ledger squeeze,
it may escape every stage-independent bound.  If there is no
exhaustive-ledger squeeze, the union of the moving supports may cover
all of \(\mathbb R^3\).  Any one of these records supplies the missing
compactness or domination statement; without one, moving-support data
does not promote to carrier closure.
\end{proof}

\begin{proof}[Proof of \Cref{thm:ns:no-go-published-analyticity}]
The half-log import gives
\(\rho(t)\gtrsim \log\lambdaJ/\lambdaJ\), hence
\(\sigmaJ=\lambdaJ\rho(t)/\CG\gtrsim \log\lambdaJ/\CG\).  This meets
the basic radius-window gate of \Cref{thm:ns:gevrey-dominance}; it is
not a failure of the gate.  The direct predictive-contraction route is
stronger.  The Gevrey optimization leaves a saddle correction of order
\(\sqrt{\log\lambdaJ}\) in the staged budget, and
\Cref{thm:ns:no-go-half-log-endpoint} computes that the resulting
half-log stage ratio tends to \(2^{2p}>1\).  Thus direct contraction
would require an additional analyticity-radius improvement, equivalently
a \(\sqrt{\log\lambdaJ}\)-scale gain beyond the available lane.  The
Littlewood--Paley transport rescue supplies the missing decay by moving
the budget to \(H^{-s}\), not by claiming a stronger radius.
\end{proof}

\begin{proof}[Proof of \Cref{thm:ns:no-go-raw-nonlinearity}]
The Schur complement is typed for linear maps \(L:E_{J,t}\to Y\).
A raw nonlinear observable \(F(u)\), such as a flux, pressure
interaction, vortex-stretching term, or triad, is not such a map from
the carrier to a probe codomain.  Consequently the Gram block
\(K_{FF}\) is not defined in the residual calculus, and the expression
\[
  K_{DD}-K_{DF}K_{FF}^{\dagger}K_{FD}
\]
is not a well-formed Schur complement.  Nonlinear content can still
enter lawfully in exactly the declared ways: tangent-linearize \(F\)
around the audited solution arc to obtain a linear probe; lift the
interaction to an augmented carrier where the interaction becomes a
linear coordinate; or record the nonlinear contribution as a source
entry in \(\Omsrc\).  What is excluded is only the untyped insertion of
raw nonlinear data into the native family.
\end{proof}

The main line of the argument remains in
\Cref{sec:adequacy-cascade,sec:no-go,sec:gevrey-lp}, where the
cascade is built, shortcut payments are excluded, and the Gevrey/LP
rescue pays the remaining amplitude obligation.

 \section{Lean Formalization Disclosure}
\label{app:ns-formalization}

\subsection{Trust base and methodology}
\label{subsec:paper3-trust-methodology}

The Lean artifacts cited in this appendix were checked with Lean 4.28.0.
The project-wide trust base allowlisted in
\texttt{lean/manifests/trust\_base.txt} comprises five core kernel axioms:
\texttt{propext}, \texttt{Classical.choice}, \texttt{Quot.sound},
\texttt{choice}, and \texttt{funext}.  An axiom-closure audit
(\texttt{\#print axioms}) of every NS-axis theorem cited in this appendix
shows the actual dependency set is contained in the strictly smaller
subset $\{\texttt{propext},\ \texttt{Classical.choice},\ \texttt{Quot.sound}\}$;
several no-go theorems depend on no axioms beyond the kernel default.  The
project discipline forbids \texttt{sorry}, project-local \texttt{axiom}, and
\texttt{opaque} declarations in the needles Lean corpus; the trust-base
allowlist is the manifest checker's gate when auditing theorem dependency
closures.  Coverage is recorded per row as \emph{theorem}, \emph{partial},
or \emph{definition}; these labels govern the wording below.

\subsection{Mechanization disclosure tiers}
\label{subsec:paper3-mechanization-disclosure-tiers}

Body footnotes use a three-tier disclosure discipline keyed to the
alignment column of the per-label coverage table.

\begin{table}[H]
  \centering
  \small
  \begin{tabularx}{\linewidth}{>{\raggedright\arraybackslash}p{0.25\linewidth} >{\raggedright\arraybackslash}p{0.11\linewidth} >{\raggedright\arraybackslash}X}
    \toprule
    Alignment & Tier label & Body footnote wording template \\
    \midrule
    \texttt{faithful} & T1 & Verified in Lean as \(\langle\)declaration\(\rangle\); see \Cref{subsec:paper3-per-label-coverage}. \\
    \texttt{narrowed-surrogate / partial} & T2 & Mechanized in Lean as \(\langle\)declaration\(\rangle\); the Lean theorem proves \(\langle\)sub-claim\(\rangle\) under \(\langle\)elevation\(\rangle\).  The full statement is proved here.  See \Cref{subsec:paper3-per-label-coverage}. \\
    \texttt{projection-packaged} & T3 & Tracked in the Lean formalization as \(\langle\)declaration\(\rangle\): a conditional schema over \(\langle\)record fields\(\rangle\) rather than an independent derivation, which is proved here.  See \Cref{subsec:paper3-per-label-coverage}. \\
    \bottomrule
  \end{tabularx}
  \caption{Paper 3 body footnote disclosure tiers.}
  \label{tab:paper3-mechanization-disclosure-tiers}
\end{table}

This tier discipline is the canonical wording rule for body footnotes:
a reader can cross-check each theorem footnote against the alignment
column of the coverage table below.

\subsection{Standing narrowings and elevated hypotheses}
\label{subsec:paper3-standing-narrowings}

The NS axis runs on the rational-finite-matrix infrastructure shared with
the Main and Xi axes; that infrastructure does not include a real-analysis
stack (limits, $\inf$, $\sqrt{\cdot}$, $\exp$, $\log$, asymptotic
$\lesssim$).  The Gevrey, half-log, LP-transport, and BG-staging theorems
therefore carry their analytic content as \emph{hypotheses} (closed-form
identities, optimization-step implications, ratio-contraction
certificates, summability witnesses) supplied by the paper-grounded
PDE-import structures of \texttt{SixBirdsNeedles.NS.PDEImports}.  Lean
mechanizes the structural composition (implication chains, record-field
unpacking, modus tollens for the no-gos); the paper supplies the
substantive estimates these compositions consume.

The six PDE imports (Beale--Kato--Majda, Bradshaw--Grujić, Kato local
well-posedness, Foias--Temam / Herbst--Skibsted / Wang half-log
analyticity, Littlewood--Paley / Bony--Chemin--Danchin calculus, and the
Sobolev/Besov embedding chain) are each realized as a Lean
\texttt{structure} whose constructor consumes the named hypotheses and
exposes the conclusion as fields.  No project-local \texttt{axiom},
\texttt{opaque}, or \texttt{sorry} appears in any of these structures;
the boundary policy is \texttt{external\_pde\_import} per
\texttt{paper/notation\_and\_terminology.md} §8.

Five of the nine closure records ($R_1, R_2, R_3, R_7, R_8, R_9$) are
realized as proof-carrying records whose component propositions are
schema slots filled by downstream consumers; three ($R_4, R_5, R_6$)
carry Lean-typed data (the source-admission ledger, the
frequency-doubling equation, the LP-transport stage operator) that
constrains the rest of the calculus.  The bundle
\texttt{ClosureR1R9} contains exactly these nine record fields and
does \emph{not} carry the conclusion $\Tmax = \infty$, which is the
downstream content of T10.

The framework-deliverable regularity theorem T10 is
\emph{projection-packaged}: it takes a \texttt{ClosureR1R9} bundle, the
smooth-data hypotheses on $u_0$, and a typed-context deliverability
bridge (an implication consuming every R-record field plus the
smooth-data data and concluding $\Tmax(u_0) = \infty$), and produces
the conclusion by unpacking the bundle and applying the bridge.  The
substantive Navier--Stokes content lives in the bridge hypothesis; Lean
audits the dependency chain.

\subsection{Per-label coverage tables}
\label{subsec:paper3-per-label-coverage}

Tables~\ref{tab:paper3-lean-coverage-defs} and~\ref{tab:paper3-lean-coverage-thms}
record the per-row Lean coverage, alignment, and short disclosure note for each
statements-of-record row.  After the 2026-05-24 NS-AOR bridge restructure,
the first table covers 15 definitions (the original 14 plus the new
\path{def:ns:aor-instance-object}) and the second covers 22 theorems
(the original 16 plus 6 NS-AORInstance theorems, one of which is partial).

\begin{table}[H]
  \centering
  \scriptsize
  \setlength{\tabcolsep}{3pt}
  \renewcommand{\arraystretch}{1.0}
  \begin{tabularx}{\linewidth}{>{\raggedright\arraybackslash}p{0.36\linewidth} >{\raggedright\arraybackslash}p{0.11\linewidth} >{\raggedright\arraybackslash}X}
    \toprule
    Paper label & Alignment & Disclosure note \\
    \midrule
    \path{def:ns:bg-readout} & faithful & Realized in Lean as a subtype with the right-support constraint $\DBG = \DBG \cdot P_{J(t)}$. \\
    \path{def:ns:physical-native-probes} & faithful & Realized in Lean as a subtype with LP-shell / Bessel-potential tagging and the energy+enstrophy currency-recovery condition. \\
    \path{def:ns:bg-package} & faithful & Realized in Lean as a subtype with admissibility, divergence-free preservation, exact energy/response, and the doubling equation. \\
    \path{def:ns:null-mode-legal-ns} & faithful & Realized in Lean as the conjunction of \texttt{NullLegal} for $\Lphys$ and $\DBG$. \\
    \path{def:ns:bkm-import} & external PDE & Realized in Lean as a proof-carrying structure for the BKM continuation criterion. \\
    \path{def:ns:bg-import} & external PDE & Realized in Lean as a proof-carrying structure for the Bradshaw--Grujić window-bound criterion. \\
    \path{def:ns:kato-import} & external PDE & Realized in Lean as a proof-carrying structure for Kato local well-posedness on $H^{\gamma_0}$ with $\gamma_0>5/2$. \\
    \path{def:ns:wang-halflog-import} & external PDE & Realized in Lean as a proof-carrying structure for the half-log analyticity radius bound. \\
    \path{def:ns:lp-bcd-import} & external PDE & Realized in Lean as a proof-carrying structure for the Littlewood--Paley / Bony--Chemin--Danchin calculus. \\
    \path{def:ns:sobolev-besov-embedding-import} & external PDE & Realized in Lean as a proof-carrying structure for the $H^{\gamma_0}\hookrightarrow\dot B^{-\eps}_{\infty,\infty}$ embedding chain. \\
    \path{def:ns:source-admission-ledger} & faithful & Realized in Lean as a structure with the five positive components, the decomposition equation, and PSD constraints. \\
    \path{def:ns:gevrey-probes} & faithful & Realized in Lean as a structure with the normalized derivative formula. \\
    \path{def:ns:lp-transport} & faithful & Realized in Lean as a structure with the stage map, exponent condition $s\ge q-\eps$, and norm-bound constant. \\
    \path{def:ns:closure-r1-r9} & faithful & Realized in Lean as nine separate proof-carrying structures plus the \texttt{ClosureR1R9} bundle; the bundle does not carry the $\Tmax=\infty$ conclusion. \\
    \path{def:ns:aor-instance-object} & faithful & Realized in Lean as the fixed-scope PreSB / RefPreSB object that packages the BG-stage chain and the audit registries for the AOR-instance reading; see \path{SixBirdsNeedles.NS.AORInstance.aorInstanceObject}. \\
    \bottomrule
  \end{tabularx}
  \caption{Definition coverage (post-2026-05-24 NS-AOR bridge restructure: 14 original NS definitions plus the new \texttt{def:ns:aor-instance-object}; all 15 rows have status \texttt{definition}).}
  \label{tab:paper3-lean-coverage-defs}
\end{table}

\begin{table}[H]
  \centering
  \scriptsize
  \setlength{\tabcolsep}{3pt}
  \renewcommand{\arraystretch}{1.0}
  \begin{tabularx}{\linewidth}{>{\raggedright\arraybackslash}p{0.34\linewidth} >{\raggedright\arraybackslash}p{0.09\linewidth} >{\raggedright\arraybackslash}p{0.13\linewidth} >{\raggedright\arraybackslash}X}
    \toprule
    Paper label & Lean coverage & Alignment & Disclosure note \\
    \midrule
    \path{thm:ns:criterion-native-equality} & theorem & faithful & Lean proves the trivial algebraic equality $\Xi_C(\DBG\mid\DBG)=0$ from the Moore--Penrose support identity. \\
    \path{thm:ns:no-go-criterion-native} & partial & narrowed-surrogate & Lean witnesses the no-go via an explicit $2$--$1$--$1$ finite counterexample where the criterion-native residual vanishes but the physical-native residual is nonzero. \\
    \path{thm:ns:source-admission-recurrence} & theorem & faithful & Lean mechanizes the unbundling of the accepted Loewner recurrence into the five-component form. \\
    \path{thm:ns:no-go-scalar-energy} & theorem & faithful & Lean mechanizes the modus-tollens chain under the linear-family hypothesis. \\
    \path{thm:ns:no-go-packing-amplitude} & theorem & faithful & Lean mechanizes the case-split modus tollens under the threshold / bridge / tail disjunction. \\
    \path{thm:ns:no-go-sparse-shadow} & theorem & faithful & Lean mechanizes the case-split modus tollens under the sector-containment / paid-mismatch disjunction. \\
    \path{thm:ns:no-go-marker-only} & theorem & faithful & Lean mechanizes the case-split modus tollens under the new-native-probe / record-grade-bound disjunction. \\
    \path{thm:ns:no-go-metric-shadow} & theorem & faithful & Lean mechanizes the modus tollens under the normalized-dominance hypothesis. \\
    \path{thm:ns:no-go-moving-window} & theorem & faithful & Lean mechanizes the three-branch case-split for finite-termination / fixed-ledger / exhaustive-ledger promotion. \\
    \path{thm:ns:no-go-half-log-endpoint} & theorem & faithful & Lean mechanizes the modus tollens under the stage-ratio-above-one hypothesis. \\
    \path{thm:ns:no-go-published-analyticity} & theorem & faithful & Lean mechanizes the three-branch case-split for LP-transport / full-log / alternative-radius rescue. \\
    \path{thm:ns:no-go-raw-nonlinearity} & theorem & faithful & Lean mechanizes the three-branch case-split for tangent-linearization / interaction-lifting / source-admission entry. \\
    \path{thm:ns:gevrey-dominance} & theorem & faithful & Lean mechanizes the closed-form / optimization / radius-window-gate composition chain. \\
    \path{thm:ns:lp-transport-rescue} & theorem & faithful & Lean mechanizes the transport-record $\to$ ratio-contraction $\to$ summability composition chain. \\
    \path{thm:ns:smooth-data-half-log-transfer} & theorem & faithful & Lean mechanizes the Sobolev/Besov embedding $\to$ Wang half-log radius composition under the all-six-channels audit. \\
    \path{thm:ns:framework-deliverable-regularity} & theorem & projection-packaged & Lean audits the dependency chain: the bundle unpacks into a deliverability-bridge hypothesis whose conclusion is $\Tmax(u_0)=\infty$. \\
    \bottomrule
  \end{tabularx}
  \caption{Theorem coverage, original NS rows (16 theorems pre-restructure).}
  \label{tab:paper3-lean-coverage-thms}
\end{table}

\begin{table}[H]
  \centering
  \scriptsize
  \setlength{\tabcolsep}{3pt}
  \renewcommand{\arraystretch}{1.0}
  \begin{tabularx}{\linewidth}{>{\raggedright\arraybackslash}p{0.34\linewidth} >{\raggedright\arraybackslash}p{0.09\linewidth} >{\raggedright\arraybackslash}p{0.13\linewidth} >{\raggedright\arraybackslash}X}
    \toprule
    Paper label & Lean coverage & Alignment & Disclosure note \\
    \midrule
    \path{thm:ns:aor-mechanical-r-records} & theorem & faithful & Lean records the mechanical R1/R2/R3/R5/R9 by-construction and R8 approved-other discharges as the mechanical components of the AOR-membership witness; the body proof carries the framing. \\
    \path{thm:ns:aor-r4-discharge} & theorem & faithful & Lean records the R4 source-admission discharge with primary type $\Delta_{\mathrm{source}}$, consuming the carrier-space quotient-transport record; the body proof presents the membership-component reading. \\
    \path{thm:ns:aor-r6-discharge} & theorem & faithful & Lean records the R6 LP-transport discharge with primary type $\Delta_{\mathrm{transport}}$ and the contractive-tail / summable secondary statuses. \\
    \path{thm:ns:aor-r7-discharge} & theorem & faithful & Lean records the R7 half-log discharge with primary type $\Delta_{\mathrm{scale}}$ and the asymptotic-budgeted basic-gate / out-of-scope-horizon secondary statuses. \\
    \path{thm:ns:aor-quotient-transport-record} & theorem & faithful & Lean records the $R_{4q}$ quotient-transport record using the carrier-space audit forms and the kernel-preservation field from \texttt{R4SourceAdmission}; the smooth-kernel-zero clause is a PSD narrowing pending a future positive-definite strengthening. \\
    \path{thm:ns:in-refstableaor} & partial & projection-packaged & Lean projects the joint discharges through the carrier-attached \texttt{AORCarrierWithDischarge} wrapper; the per-atom assignment dispatches by primary residual type and routes through generic \path{reachableFromZero}/\path{residualRecordOfAtom_eventuallyEmpty} lemmas because the R-record acceptance fields are \texttt{Prop}-valued.  Tighter content-dependence on R-record values would require making the R-record acceptance fields data-bearing.  See \Cref{subsec:paper3-in-refstableaor-narrowing} for the partial-demotion narrowing. \\
    \bottomrule
  \end{tabularx}
  \caption{Theorem coverage, NS-AORInstance additions (6 theorems added 2026-05-24 by the NS-AOR bridge restructure; one partial demotion on the joint refinement-stable AOR statement).}
  \label{tab:paper3-lean-coverage-thms-aor-instance}
\end{table}

\subsection{Lean declaration map}
\label{subsec:paper3-lean-declaration-map}

Table~\ref{tab:paper3-lean-declaration-map} maps each paper label to
its full Lean declaration path; the per-row coverage notes above use
short prose disclosures and refer to this table for the canonical
identifiers.

\begin{table}[H]
  \centering
  \scriptsize
  \setlength{\tabcolsep}{3pt}
  \renewcommand{\arraystretch}{1.0}
  \begin{tabularx}{\linewidth}{>{\raggedright\arraybackslash}p{0.36\linewidth} >{\raggedright\arraybackslash}X}
    \toprule
    Paper label & Lean declaration path \\
    \midrule
    \path{def:ns:bg-readout} & \path{SixBirdsNeedles.NS.Carrier.bgReadout} \\
    \path{def:ns:physical-native-probes} & \path{SixBirdsNeedles.NS.Carrier.physicalNativeProbes} \\
    \path{def:ns:bg-package} & \path{SixBirdsNeedles.NS.Carrier.bgPackage} \\
    \path{def:ns:null-mode-legal-ns} & \path{SixBirdsNeedles.NS.Carrier.nullModeLegalNS} \\
    \path{def:ns:bkm-import} & \path{SixBirdsNeedles.NS.PDEImports.bkmImport} \\
    \path{def:ns:bg-import} & \path{SixBirdsNeedles.NS.PDEImports.bgImport} \\
    \path{def:ns:kato-import} & \path{SixBirdsNeedles.NS.PDEImports.katoImport} \\
    \path{def:ns:wang-halflog-import} & \path{SixBirdsNeedles.NS.PDEImports.wangHalfLogImport} \\
    \path{def:ns:lp-bcd-import} & \path{SixBirdsNeedles.NS.PDEImports.lpBcdImport} \\
    \path{def:ns:sobolev-besov-embedding-import} & \path{SixBirdsNeedles.NS.PDEImports.sobolevBesovEmbeddingImport} \\
    \path{thm:ns:criterion-native-equality} & \path{SixBirdsNeedles.NS.CriterionNativeGap.criterionNativeEquality} \\
    \path{thm:ns:no-go-criterion-native} & \path{SixBirdsNeedles.NS.CriterionNativeGap.noGoCriterionNative} \\
    \path{def:ns:source-admission-ledger} & \path{SixBirdsNeedles.NS.Cascade.sourceAdmissionLedger} \\
    \path{thm:ns:source-admission-recurrence} & \path{SixBirdsNeedles.NS.Cascade.sourceAdmissionRecurrence} \\
    \path{thm:ns:no-go-scalar-energy} & \path{SixBirdsNeedles.NS.NoGo.noGoScalarEnergy} \\
    \path{thm:ns:no-go-packing-amplitude} & \path{SixBirdsNeedles.NS.NoGo.noGoPackingAmplitude} \\
    \path{thm:ns:no-go-sparse-shadow} & \path{SixBirdsNeedles.NS.NoGo.noGoSparseShadow} \\
    \path{thm:ns:no-go-marker-only} & \path{SixBirdsNeedles.NS.NoGo.noGoMarkerOnly} \\
    \path{thm:ns:no-go-metric-shadow} & \path{SixBirdsNeedles.NS.NoGo.noGoMetricShadow} \\
    \path{thm:ns:no-go-moving-window} & \path{SixBirdsNeedles.NS.NoGo.noGoMovingWindow} \\
    \path{thm:ns:no-go-half-log-endpoint} & \path{SixBirdsNeedles.NS.NoGo.noGoHalfLogEndpoint} \\
    \path{thm:ns:no-go-published-analyticity} & \path{SixBirdsNeedles.NS.NoGo.noGoPublishedAnalyticity} \\
    \path{thm:ns:no-go-raw-nonlinearity} & \path{SixBirdsNeedles.NS.NoGo.noGoRawNonlinearity} \\
    \path{def:ns:gevrey-probes} & \path{SixBirdsNeedles.NS.GevreyLP.gevreyProbes} \\
    \path{def:ns:lp-transport} & \path{SixBirdsNeedles.NS.GevreyLP.lpTransport} \\
    \path{thm:ns:gevrey-dominance} & \path{SixBirdsNeedles.NS.GevreyLP.gevreyDominance} \\
    \path{thm:ns:lp-transport-rescue} & \path{SixBirdsNeedles.NS.GevreyLP.lpTransportRescue} \\
    \path{thm:ns:smooth-data-half-log-transfer} & \path{SixBirdsNeedles.NS.GevreyLP.smoothDataHalfLogTransfer} \\
    \path{def:ns:closure-r1-r9} & \path{SixBirdsNeedles.NS.ClosureRecords.ClosureR1R9} \\
    \path{thm:ns:framework-deliverable-regularity} & \path{SixBirdsNeedles.NS.MainTheorem.frameworkDeliverableRegularity} \\
    \path{def:ns:aor-instance-object} & \path{SixBirdsNeedles.NS.AORInstance.aorInstanceObject} \\
    \path{thm:ns:aor-mechanical-r-records} & \path{SixBirdsNeedles.NS.AORInstance.aorMechanicalRRecords} \\
    \path{thm:ns:aor-r4-discharge} & \path{SixBirdsNeedles.NS.AORInstance.aorR4Discharge} \\
    \path{thm:ns:aor-r6-discharge} & \path{SixBirdsNeedles.NS.AORInstance.aorR6Discharge} \\
    \path{thm:ns:aor-r7-discharge} & \path{SixBirdsNeedles.NS.AORInstance.aorR7Discharge} \\
    \path{thm:ns:aor-quotient-transport-record} & \path{SixBirdsNeedles.NS.AORInstance.aorQuotientTransportRecord} \\
    \path{thm:ns:in-refstableaor} & \path{SixBirdsNeedles.NS.AORInstance.inRefStableAOR} \\
    \bottomrule
  \end{tabularx}
  \caption{Mapping from paper labels to Lean declaration paths (post-2026-05-24 NS-AOR bridge restructure: 37 rows including the 7 new NS-AORInstance entries).}
  \label{tab:paper3-lean-declaration-map}
\end{table}

\subsection{Partial-demotion narrowing for \texorpdfstring{\texttt{thm:ns:in-refstableaor}}{thm:ns:in-refstableaor}}
\label{subsec:paper3-in-refstableaor-narrowing}

Of the seven new coverage rows recorded above for the
NS-as-AOR-instance reading (see \Cref{sec:ns-aor-instance}), six
discharge with \texttt{status = "theorem"} and one,
\path{thm:ns:in-refstableaor}, carries \texttt{status = "partial"}.
The partial status is disclosed here for transparency.

The Lean proof in
\path{SixBirdsNeedles.NS.AORInstance.inRefStableAOR} case-analyzes
by residual primary type and brings the R4 / R6 / R7 discharge
fields into lexical scope.  Each case then closes through the
generic lemmas \path{reachableFromZero_zero},
\path{reachableFromZero_budgeted}, and
\path{residualRecordOfAtom_eventuallyEmpty}.  These generic closures
are what make the proof projection-packaged rather than content-bearing:
the per-atom assignment dispatches by primary type (and routes the
non-standalone branches through the cross-cascade), but it does not
pattern-match against the R-record data fields.

The deferral lives in
\path{SixBirdsNeedles.NS.ClosureRecords}.  The R-record acceptance
fields (\path{recurrenceAccepted}, \path{typedTransportAccepted},
\path{halfLogAnalyticityAccepted}, \path{bgCompatibleStages},
\path{smoothDataSubcarrier}, etc.) are \texttt{Prop}-valued
predicates rather than data-bearing structures.  Closing the gap
structurally would require an NS-side surgery to make those fields
data-bearing (e.g.\ promoting them to \texttt{Bool}-valued or to
explicit data), with downstream consequences for every consumer of
\texttt{R4SourceAdmission}, \texttt{R6LPTransport}, and
\texttt{R7HalfLogAnalyticity}.  Such a strengthening is outside the
scope of the present manuscript; the \texttt{partial} status is the
honest disclosure that the AOR-instance reading currently rests on
\texttt{Prop}-valued audit-record witnesses.

\subsection{Carrier-space field additions to \texorpdfstring{\texttt{R4SourceAdmission}}{R4SourceAdmission}}
\label{subsec:paper3-r4-carrier-space-fields}

The structure \path{SixBirdsNeedles.NS.ClosureRecords.R4SourceAdmission}
carries carrier-space fields used by the
\path{aorQuotientTransportRecord} theorem in
\path{SixBirdsNeedles.NS.AORInstance}: \path{e}, \path{eNext}
(carrier-space dimensions), \path{C_J_t}, \path{C_J_t_h} (audit
forms), \path{audit_psd_t}, \path{audit_psd_t_h} (positive-semidefinite
proofs), \path{carrierTransport} (the audited transport map
\(\texttt{Mat eNext e}\)), and \path{kernel_preserved} (the predicate
\(T_h(\Ker C_{J,t}) \subseteq \Ker C_{J,t+h}\)).

These fields were originally added to support a bridge theorem
hosted in the companion meta-theory manuscript
\citet{Tsiokos2026AOR}.  Both producer (\path{R4SourceAdmission})
and consumer (\path{aorQuotientTransportRecord}) now live in the
NS-axis Lean development, and the field addition is recorded here on
the NS side.
 
\clearpage
\bibliographystyle{plainnat}
\bibliography{references}

\end{document}
